\chapter{A transcendental Nakamaye theorem}
\label{chap:transcendental-nakamaye}
\providecommand{\NkRes}[1]{\operatorname{Res}_{#1}}

The singularity theory of the preceding chapters has a concrete geometric
application: it identifies the non-K\"ahler locus of a pseudoeffective
class through restricted volumes. We present the full argument of the
manuscript of Vu and Xia supplied for this book \cite{VXNull}, using
\cite{SGPT} for its trace and singularity-space theory. The ambient
manifold in the main theorem is K\"ahler. The quantitative regularization
lemma below is proved for a closed background on an arbitrary compact
Hermitian manifold; the numerical and trace arguments retain their
K\"ahler hypotheses.

\section{Statement and strategy}

Let $X$ be a connected compact K\"ahler manifold of positive dimension,
and let $\alpha$ be a pseudoeffective ordinary real Bott--Chern class.
Use the numerical restricted volume of
\cref{eq:bc-numerical-restricted-volume}, with its smooth-remainder
convention. Define
\begin{equation}\label{nk:eq:null-locus}
 \operatorname{Null}(\alpha)
 =\bigcup_{\substack{V\subset X\ \mathrm{irreducible},\ \dim V>0\\
          \langle\alpha^{\dim V}\rangle_{X\mid V}=0}}V.
\end{equation}

\begin{theorem}\label{nk:thm:main}
If $V$ is a positive-dimensional irreducible component of
$E_{\mathrm{nK}}(\alpha)$, then
\[
 \langle\alpha^{\dim V}\rangle_{X\mid V}=0.
\]
Consequently,
$E_{\mathrm{nK}}(\alpha)=\operatorname{Null}(\alpha)$.
\end{theorem}

The vanishing assertion concerns irreducible components. Equality with
the union in \eqref{nk:eq:null-locus} does not assert vanishing on every
subvariety contained in $E_{\mathrm{nK}}(\alpha)$.

Suppose that a component has positive restricted volume. After resolving
it, the trace formula expresses this volume as a positive mass on a
smooth submanifold. A rooftop envelope retains a fixed positive part of
that mass and supplies a strictly positive current. Multiplier ideals
then make its singularities comparable with those of an ambient current.
On a suitable modification, a pointwise comparison on a collar allows
these currents to be glued. The resulting current is smooth at a point
of the strict transform, contradicting the transformation rule for the
non-K\"ahler locus.

The proof uses a coarse approximation level $m$ and a much finer level
$M$. The constants controlling positivity are fixed before $M$ is
chosen. This order is recorded explicitly in
\cref{nk:sec:fixed-data,nk:sec:fixed-principalization}; it prevents the
positivity gain from being lost when singularities are compared.

\section{Preliminaries}\label{nk:sec:prelim}

Throughout, $X$ denotes a connected compact K\"ahler manifold, and $\omega$ denotes a K\"ahler form on $X$.
All complex submanifolds are understood to be closed; in particular,
submanifolds of compact manifolds are compact. We
write $\ddc=(2\pi)^{-1}\mathrm{i}\partial\bar\partial$, so that
$\ddc\log|z|^2$ is the current of integration along the origin in
$\mathbf C$. A \emph{modification} is a proper bimeromorphic morphism between
compact complex manifolds which is a finite composition of blowups with
smooth centers. A \emph{prime divisor over $X$} is a prime divisor
$E$ on some modification $\pi\colon Y\to X$. For a quasi-plurisubharmonic
(quasi-psh) function $\varphi$ on $X$ we write $\nu(\varphi,E)$ for the
generic Lelong number of $\pi^*\varphi$ along $E$. If
$Z\subset X$ is an irreducible analytic subset, $\nu(\varphi,Z)$ denotes
the generic Lelong number of $\varphi$ along $Z$. If $D$ is a divisor
on a modification and $f$ is a further modification, $f^*D$ denotes the
total transform of $D$.

\subsection{Quasi-psh functions, multiplier ideals and analytic
singularities}\label{nk:sec:qpsh}

For a smooth closed real $(1,1)$-form $\theta$ on $X$, we denote by
$\PSH(X,\theta)$ the set of $\theta$-psh functions
$\varphi\colon X\to[-\infty,\infty)$, that is, quasi-psh functions in
their canonical upper semicontinuous representatives, not identically
$-\infty$, with $\theta_\varphi\coloneqq \theta+\ddc\varphi\ge0$.
Equivalently, if $\theta=\ddc h$ locally, then $h+\varphi$ is psh.
For quasi-psh functions
$\varphi,\psi$ we write $\varphi\preceq\psi$ if
$\varphi\le\psi+C$ for a constant $C$, and we say that $\varphi$ is
\emph{more singular} than $\psi$; we write $\varphi\sim\psi$ if
$\varphi\preceq\psi$ and $\psi\preceq\varphi$.

For a quasi-psh function $\varphi$ and $\lambda>0$, the
\emph{multiplier ideal} $\mathcal I(\lambda\varphi)\subset\mathcal O_X$, introduced
by Nadel \cite{Nad90}, is the sheaf of
germs of holomorphic functions $f$ with
$|f|^2\mathrm{e}^{-\lambda\varphi}\in L^1_{\mathrm{loc}}$. It is coherent and
integrally closed. By the subadditivity theorem \cite{DEL00},
\begin{equation}\label{nk:eq:subadditivity}
  \mathcal I\bigl((\lambda+\lambda')\varphi\bigr)
  \subset\mathcal I(\lambda\varphi)\,\mathcal I(\lambda'\varphi),
  \qquad \lambda,\lambda'>0.
\end{equation}
For a coherent ideal $\mathfrak a$ with local generators $f_1,\dots,f_N$, we
set
\[
  \psi_{\mathfrak a}\coloneqq \log\sum_{j=1}^N|f_j|^2;
\]
a different choice of generators changes $\psi_\mathfrak a$ by a locally bounded
function; we write $O(1)$ for a locally bounded function. If
$\mathfrak b\subset\mathfrak a$ are coherent ideals, then
$\psi_{\mathfrak b}\le\psi_{\mathfrak a}+O(1)$, because local generators
of $\mathfrak b$ are combinations of local generators of $\mathfrak a$
with holomorphic coefficients; moreover
$\psi_{\mathfrak a\mathfrak b}=\psi_{\mathfrak a}+\psi_{\mathfrak b}+O(1)$,
and $\psi_{\mathfrak a}\circ\pi=\psi_{\mathfrak a\mathcal O_Y}+O(1)$ for
every modification $\pi\colon Y\to X$, because the pullbacks of local
generators of $\mathfrak a$ generate $\mathfrak a\mathcal O_Y$. If
$Y\subset X$ is a closed submanifold, the
\emph{restriction} of $\mathfrak a$ to $Y$ is the image ideal
$\NkRes{Y}\mathfrak a\coloneqq \operatorname{im}(\mathfrak a\to\mathcal O_Y)\subset\mathcal O_Y$
\cite[Definition~1.4.5]{SGPT}; it is generated by the restrictions of
local generators of $\mathfrak a$, and it is not to be confused with
$\mathfrak a\otimes\mathcal O_Y$. For an effective divisor $D$ on $X$ we write
$\psi_D\coloneqq \log|s_D|_{h_D}^2$, where $s_D$ is the canonical section of
$\mathcal O_X(D)$ and $h_D$ is a smooth Hermitian metric on $\mathcal O_X(D)$; by the
Poincar\'e--Lelong formula, $\ddc\psi_D=[D]-\Theta_D$, where $\Theta_D$
denotes the Chern form of $h_D$. The same notation is used for effective
$\mathbf Q$-divisors, after clearing denominators. If $D_1\ge D_2$ are
effective $\mathbf Q$-divisors, then
$\psi_{D_1}-\psi_{D_2}=\psi_{D_1-D_2}+O(1)$
is bounded above on $X$, and if $\mathfrak a\mathcal O_Y=\mathcal O_Y(-D)$ for
a modification $\pi\colon Y\to X$, then
$\psi_{\mathfrak a}\circ\pi=\psi_D+O(1)$.
When an exact identity between divisorial potentials or their Chern forms
is used, the metrics are chosen compatibly under tensor products and
pullback; arbitrary choices differ by smooth bounded terms.

A quasi-psh function $\varphi$ has \emph{analytic singularities of type
$(c,\mathfrak a)$}, where $c\in\mathbf Q_{>0}$ and $\mathfrak a\subset\mathcal O_X$ is a coherent
ideal, if locally $\varphi=c\,\psi_\mathfrak a+R$ with $R$ bounded
\cite[Definition~1.6.1]{SGPT}; the remainder $R$ is merely bounded.
(The definition of numerical restricted volumes in
\cref{nk:sec:restricted-volumes} uses the stronger convention of
\cite[Section~2]{CT22Restricted}, with smooth remainder, which we recall there.)
If $\varphi$ has analytic singularities of type
$(c,\mathfrak a)$ and $\pi\colon Y\to X$ is a modification such that
$\mathfrak a\mathcal O_Y=\mathcal O_Y(-D)$ is invertible (a \emph{principalization} of
$\mathfrak a$), then $\pi^*\varphi=c\,\psi_D+R'$ with $R'$ bounded, and we
say that $\pi^*\varphi$ has \emph{log singularities} along the
$\mathbf Q$-divisor $cD$. Such a $\pi$ always exists for a coherent ideal
on a compact complex manifold, and any modification dominating a
principalization is again a principalization. If $\varphi$ has analytic
singularities of type $(c,\mathfrak a)$ and $Y\subset X$ is a submanifold with
$\mathfrak a\not\subset\mathcal I_Y$, then $\varphi|_Y\not\equiv-\infty$ and
$\varphi|_Y$ has analytic singularities of type $(c,\NkRes{Y}\mathfrak a)$.

A quasi-psh function $\varphi$ with analytic singularities on a compact
manifold admits a \emph{log resolution},
that is, a modification $\pi\colon Y\to X$ such that $\pi^*\varphi$ has
log singularities along an effective $\mathbf Q$-divisor with simple normal
crossing support \cite[Theorem~1.6.1]{SGPT}.

The following consequence of log singularities will be used several
times; see \cite[Lemma~1.6.1]{SGPT}.

\begin{lemma}\label{nk:lem:residual-nef}
Let $X$ be a compact K\"ahler manifold, let $\theta$ be a smooth
closed real $(1,1)$-form on $X$, and let $\varphi\in\PSH(X,\theta)$
have log singularities along an effective $\mathbf Q$-divisor $D$.
Then the class $\{\theta\}-\{D\}$ is nef. If moreover
$\theta+\ddc\varphi$ is a K\"ahler current, then $\{\theta\}-\{D\}$
is K\"ahler.
\end{lemma}

\subsection{Non-pluripolar products, envelopes and comparison of
singularities}\label{nk:sec:npp}

Let $\theta$ be a smooth closed real $(1,1)$-form on $X$ and
$\varphi\in\PSH(X,\theta)$. The non-pluripolar product
$\theta_\varphi^n$, which goes back to Bedford--Taylor \cite{BT87} and
was developed by Guedj--Zeriahi \cite{GZ07} and
Boucksom--Eyssidieux--Guedj--Zeriahi
\cite{BEGZ10}, is a positive measure on $X$ which puts no mass on
pluripolar sets; more generally, for $\varphi_j\in\PSH(X,\theta_j)$,
$j=1,\dots,n$, the mixed non-pluripolar product
$\theta_{1,\varphi_1}\wedge\dots\wedge\theta_{n,\varphi_n}$ is always
well defined on a compact K\"ahler manifold, and it is multilinear
\cite[Proposition~2.2.1]{SGPT}. Powers of positive currents are always
understood in the non-pluripolar sense. We use the following facts.
\begin{enumerate}[label=(\alph*),leftmargin=*,itemsep=.2em,topsep=.3em]
\item\label{nk:it:monotone} (Monotonicity, due to Witt~Nystr\"om \cite{WN19}
  and Darvas--Di~Nezza--Lu \cite{DDNL18mono}; see
  \cite[Theorem~2.4.3]{SGPT} for the mixed statement.) If
  $\varphi_j\preceq\psi_j$ for all $j$, then
  $\int_X\theta_{1,\varphi_1}\wedge\dots\wedge\theta_{n,\varphi_n}
  \le\int_X\theta_{1,\psi_1}\wedge\dots\wedge\theta_{n,\psi_n}$.
\item\label{nk:it:invariance} (Invariance \cite[Proposition~2.4.2]{SGPT}.)
  If $\pi\colon Y\to X$ is a modification and $T_1,\dots,T_n$ are
  closed positive $(1,1)$-currents on $X$, then
  $\int_Y\pi^*T_1\wedge\dots\wedge\pi^*T_n=\int_XT_1\wedge\dots\wedge T_n$.
\item\label{nk:it:bigness} (\cite[Proposition~1.22]{BEGZ10}; see
  \cite[Proposition~2.4.1]{SGPT}.) If some $\varphi\in\PSH(X,\theta)$
  has positive mass $\int_X\theta_\varphi^n>0$, then the class
  $\{\theta\}$ is big.
\end{enumerate}
We write $\PSH(X,\theta)_{>0}$ for the set of $\theta$-psh functions with
positive non-pluripolar mass, and
$V_\theta\coloneqq \sup\{\varphi\in\PSH(X,\theta):\varphi\le0\}$ for the potential
with minimal singularities.

\subsubsection*{Envelopes}
For $\varphi\in\PSH(X,\theta)$, the \emph{envelope of the singularity
type} of $\varphi$ is
\[
  P_\theta[\varphi]
  \coloneqq \sup{}^*\bigl\{\psi\in\PSH(X,\theta):\psi\le0,\ \psi\preceq\varphi\bigr\}
\]
\cite[Definition~3.1.2]{SGPT}, where $\sup^*$ denotes the upper
semicontinuous regularization of the supremum; envelopes of this kind
were introduced by Ross--Witt~Nystr\"om \cite{RWN14} and studied in big
cohomology classes by Darvas--Di~Nezza--Lu
\cite{DDNL18full,DDNL18mono}.

For a function $f\colon X\to[-\infty,\infty]$, the \emph{rooftop
envelope} is
\[
  P_\theta(f)\coloneqq \sup{}^*\bigl\{\psi\in\PSH(X,\theta):
     \psi\le f\ \text{quasi-everywhere}\bigr\}
\]
\cite[Definition~2.4.1]{SGPT} (see also \cite{RWN14,DDNL18mono}),
where quasi-everywhere means outside a
pluripolar set. It is either identically $\pm\infty$ or $\theta$-psh, and
$P_\theta(f)\le f$ quasi-everywhere.

\subsubsection*{The \texorpdfstring{$P$}{P}- and
\texorpdfstring{$\mathcal I$}{I}-preorders}
Following \cite[Definition~6.1.1]{SGPT}, for quasi-psh functions
$\varphi,\psi$ we write $\varphi\preceq_P\psi$ if
$P_\theta[\varphi]\le P_\theta[\psi]$ for some (equivalently, for every)
$\theta$ with $\varphi,\psi\in\PSH(X,\theta)_{>0}$, and
$\varphi\sim_P\psi$ if both $\varphi\preceq_P\psi$ and $\psi\preceq_P\varphi$
hold. We write $\varphi\preceq_{\mathcal I}\psi$ if
$\mathcal I(\lambda\varphi)\subset\mathcal I(\lambda\psi)$ for all $\lambda>0$, and
$\varphi\sim_{\mathcal I}\psi$ if both inclusions hold. By
\cite[Proposition~6.1.2]{SGPT} (see also \cite{BFJ08}),
$\varphi\preceq_{\mathcal I}\psi$ holds if and only if
$\nu(\varphi,E)\ge\nu(\psi,E)$ for every prime divisor $E$ over
$X$. Moreover,
$\varphi\preceq\psi\Rightarrow\varphi\preceq_P\psi$ by the observation
following \cite[Definition~6.1.1]{SGPT}, while
$\varphi\preceq_P\psi\Rightarrow\varphi\preceq_{\mathcal I}\psi$ by
\cite[Corollary~6.1.1]{SGPT}. In particular, two
$\mathcal I$-equivalent functions have the same generic Lelong numbers along all
prime divisors over $X$ \cite[Proposition~3.2.1]{SGPT}. 
These preorders also apply to closed positive $(1,1)$-currents in the obvious manner.
For
$\varphi\in\PSH(X,\theta)$, the $\mathcal I$-envelope
$P_\theta[\varphi]_{\mathcal I}$, introduced by Darvas--Xia \cite{DX22},
satisfies $\varphi\sim_{\mathcal I} P_\theta[\varphi]_\mathcal I$
\cite[Proposition~3.2.6]{SGPT}.

The following theorem from \cite[Theorem~6.1.1]{SGPT} controls rooftop
envelopes of the form $P_\theta(t\varphi-(t-1)\psi)$. It supplies the
positive mass needed in \cref{nk:sec:positive-current}.

\begin{theorem}\label{nk:thm:rooftop}
Let $\varphi,\psi\in\PSH(X,\theta)_{>0}$ with $\varphi\preceq_P\psi$, and
let $t>1$ satisfy
$t^n\int_X\theta_\varphi^n>(t^n-1)\int_X\theta_\psi^n$. Then
$P_\theta\bigl(t\varphi-(t-1)\psi\bigr)$ belongs to
$\PSH(X,\theta)_{>0}$.
\end{theorem}

\subsubsection*{The pseudometric \texorpdfstring{${d_S}$}{dS}}
For a big class $\{\theta\}$, we use the pseudometric
${d_S}=d_{S,\theta}$ on $\PSH(X,\theta)$ introduced by
Darvas--Di~Nezza--Lu \cite{DDNL21}; see
\cite[Definition~6.2.1]{SGPT}. We shall only
use the following properties.
\begin{enumerate}[label=(\alph*),leftmargin=*,itemsep=.2em,topsep=.3em]
\item\label{nk:it:dS-zero} ${d_S}(\varphi,\psi)=0$ if and only if
  $\varphi\sim_P\psi$; in particular ${d_S}(\varphi,P_\theta[\varphi])=0$
  for $\varphi\in\PSH(X,\theta)_{>0}$ \cite[Proposition~6.2.2]{SGPT}.
\item\label{nk:it:dS-background} For a K\"ahler form $\omega'$, a net
  converges in $d_{S,\theta}$ if and only if it converges in
  $d_{S,\theta+\omega'}$ \cite[Corollary~6.2.9]{SGPT}; the convergence
  $\varphi_i\xrightarrow{{d_S}}\varphi$ is therefore independent of the
  background form.
\item\label{nk:it:dS-mass} If $\theta_1,\dots,\theta_n$ represent big
  classes and $\varphi_j^i\xrightarrow{{d_S}}\varphi_j$ in
  $\PSH(X,\theta_j)$ for $j=1,\dots,n$, then
  $\int_X\theta_{1,\varphi^i_1}\wedge\dots\wedge\theta_{n,\varphi^i_n}
  \to\int_X\theta_{1,\varphi_1}\wedge\dots\wedge\theta_{n,\varphi_n}$
  \cite[Theorem~6.2.1]{SGPT}.
\item\label{nk:it:dS-lelong} If $\varphi_i\xrightarrow{{d_S}}\varphi$, then
  $\nu(\varphi_i,E)\to\nu(\varphi,E)$ for every prime divisor $E$ over
  $X$ \cite[Theorem~6.2.4]{SGPT}, a result originally proved in
  \cite{Xia24}.
\end{enumerate}

\subsubsection*{Quasi-equisingular approximations}
Let $\varphi\in\PSH(X,\theta)$. A sequence $(\varphi_j)_{j\ge1}$ of
quasi-psh functions with analytic singularities is a
\emph{quasi-equisingular approximation} of $\varphi$ if it decreases to
$\varphi$, if $\varphi_j\in\PSH(X,\theta+\epsilon_j\omega)$ for a
sequence $\epsilon_j\downarrow0$, and if for all $\lambda'>\lambda>0$
one has $\mathcal I(\lambda'\varphi_j)\subset\mathcal I(\lambda\varphi)$ for $j$ large
\cite[Definition~1.6.4]{SGPT}. By Demailly's regularization theorem
\cite{Dem92,DPS01}, in the form of \cite[Theorem~1.6.2]{SGPT}, every
$\varphi\in\PSH(X,\theta)$ admits a quasi-equisingular approximation
$(\varphi_j)_j$ such that $\varphi_j$ has analytic singularities of
type $(2^{-j},\mathcal I(2^j\varphi))$. If
$\varphi\in\PSH(X,\theta)_{>0}$ and $(\varphi_j)_j$ is a decreasing
sequence of potentials with analytic singularities, with
$\varphi_j\in\PSH(X,\theta+\epsilon_j\omega)$ and $\inf_j\varphi_j=\varphi$,
then $(\varphi_j)_j$ is a quasi-equisingular approximation of
$\varphi$ if and only if $\varphi_j\xrightarrow{{d_S}}P_\theta[\varphi]_\mathcal I$
\cite[Corollary~7.1.2]{SGPT}; this circle of ideas goes back to
\cite{DX22,DX24b}.

\subsection{Non-K\"ahler locus and non-nef locus}\label{nk:sec:nK}

For a pseudoeffective class $\alpha$, the \emph{non-nef locus} is
\[
  {E_{\mathrm{nn}}}(\alpha)\coloneqq \bigcup_{\epsilon>0}{E_{\mathrm{nK}}}(\alpha+\epsilon\{\omega\}).
\]
This definition is independent of the choice of $\omega$
\cite[Section~2]{CT22Restricted}. If $\alpha$ is big and $\theta\in\alpha$ is a
smooth representative, it is equivalently given by
${E_{\mathrm{nn}}}(\alpha)=\{x\in X:\nu(V_\theta,x)>0\}$; it is contained in
${E_{\mathrm{nK}}}(\alpha)$ \cite[Definitions~1.7.6 and~1.7.7]{SGPT}. The generic
Lelong number $\nu(V_\theta,Z)$ of $V_\theta$ along an irreducible
analytic subset $Z$ vanishes if and only if $Z\not\subset{E_{\mathrm{nn}}}(\alpha)$,
since the set $\{x\in Z:\nu(V_\theta,x)>0\}$ is a countable union of
proper analytic subsets of $Z$ when $\nu(V_\theta,Z)=0$ (Siu's
theorem \cite{Siu74}). For big $\alpha$, we shall use the transformation rules for these loci under a
modification $\mu\colon\widetilde X\to X$ established in
\cite[Lemmas~4.1 and~4.3]{CT22Restricted}:
\begin{equation}\label{nk:eq:transformation-rules}
  {E_{\mathrm{nn}}}(\mu^*\alpha)=\mu^{-1}\bigl({E_{\mathrm{nn}}}(\alpha)\bigr),\qquad
  {E_{\mathrm{nK}}}(\mu^*\alpha)=\mu^{-1}\bigl({E_{\mathrm{nK}}}(\alpha)\bigr)\cup\operatorname{Exc}(\mu),
\end{equation}
where $\operatorname{Exc}(\mu)$ is the exceptional locus of $\mu$. Finally, if $\alpha$ is big, then by
\cite[Theorem~3.17]{Bou04Zariski} (see also \cite[Theorem~1.7.1]{SGPT}), there
exists a K\"ahler current $R\in\alpha$ with analytic singularities whose
polar set is exactly ${E_{\mathrm{nK}}}(\alpha)$.

\subsection{Numerical restricted volumes and the trace
operator}\label{nk:sec:restricted-volumes}

Let $\alpha$ be a pseudoeffective class, let $\theta\in\alpha$ be a
smooth representative, and let $V\subset X$ be an irreducible analytic
subset of dimension $k>0$. Following \cite[Section~2]{CT22Restricted}, the
\emph{numerical restricted volume} of $\alpha$ along $V$ is defined by
$\langle\alpha^k\rangle_{X\mid V}\coloneqq 0$ if $V\subset{E_{\mathrm{nn}}}(\alpha)$, and
otherwise by
\begin{equation}\label{nk:eq:def-restricted-volume}
  \langle\alpha^k\rangle_{X\mid V}
  \coloneqq \lim_{\epsilon\downarrow0}\ \sup_{T}
     \int_{V_{\mathrm{reg}}}
     \bigl((T+\epsilon\omega)|_{V_{\mathrm{reg}}}\bigr)_{\mathrm{ac}}^k,
\end{equation}
where $T$ runs over the closed real $(1,1)$-currents in $\alpha$
which do not contain $V$ in their polar sets, satisfy
$T\ge-\epsilon\omega$, and have analytic singularities in the strong
sense of \cite[Section~2]{CT22Restricted}: the local potentials have the form
$c\log\sum_j|f_j|^2+g$ with $c>0$ real and $g$ smooth. The
subscript $\mathrm{ac}$ denotes the absolutely continuous part. For
currents of this admissible family the integrand agrees with the
non-pluripolar product, both being the extension by zero of the smooth
product off the polar set \cite[Remark~2.2]{CT22Restricted}; we do not use, nor
assert, this identification for potentials with merely bounded
remainder. For $V=X$ one recovers the volume,
$\langle\alpha^n\rangle_{X\mid X}=\vol(\alpha)$, which vanishes if and
only if $\alpha$ is not big \cite{Bou02b}.

We shall use the following properties from \cite{CT22Restricted}.
\begin{enumerate}[label=(\alph*),leftmargin=*,itemsep=.2em,topsep=.3em]
\item\label{nk:it:rv-positive} If $V\not\subset{E_{\mathrm{nK}}}(\alpha)$, then
  $\langle\alpha^k\rangle_{X\mid V}>0$. Indeed, by
  \cite[Lemma~2.7]{CT22Restricted} the numerical restricted volume is then the
  supremum of
  $\int_{V_{\mathrm{reg}}}(T|_{V_{\mathrm{reg}}})_{\mathrm{ac}}^k$
  over K\"ahler currents $T\in\alpha$ whose local potentials have
  the form above with $g$ smooth and whose polar sets do not contain
  $V$. Such a current exists when $V\not\subset{E_{\mathrm{nK}}}(\alpha)$,
  as recalled in the proof of that lemma. If $T\ge\delta\omega$,
  its restricted integral is at least $\delta^k\int_V\omega^k>0$.
\item\label{nk:it:rv-birational} If $\mu\colon\widetilde X\to X$ is a
  modification with $V\not\subset\mu(\operatorname{Exc}(\mu))$, if
  $\widetilde V$ denotes the strict transform of $V$, and if
  $\alpha$ is big, then
  $\langle(\mu^*\alpha)^k\rangle_{\widetilde X\mid\widetilde V}
  =\langle\alpha^k\rangle_{X\mid V}$ \cite[Lemma~4.2]{CT22Restricted}.
\end{enumerate}

\subsubsection*{The trace operator}
Let $V\subset X$ be a closed connected submanifold of dimension $k$, and let
$\varphi$ be a quasi-psh function on $X$ with $\nu(\varphi,V)=0$. If
$(\varphi_j)_j$ is a quasi-equisingular approximation of $\varphi$, then
$\varphi_j|_V\not\equiv-\infty$ for all $j$, and the restrictions
$\varphi_j|_V$ converge in ${d_S}$, with respect to any sufficiently
positive K\"ahler form on $V$, to a quasi-psh function on $V$ which is
independent of the approximation up to $P$-equivalence
\cite[Proposition~8.1.1, Definition~8.1.1, Proposition~8.1.2]{SGPT}. Any
such limit is called a \emph{trace} of $\varphi$ along $V$ and is denoted
by $\operatorname{Tr}_V(\varphi)$. If $\theta$ is a smooth closed real
$(1,1)$-form with $\varphi\in\PSH(X,\theta)$, and if a trace of
$\varphi$ can be chosen in $\PSH(V,\theta|_V)_{>0}$, one defines
\[
  \operatorname{Tr}^\theta_V(\varphi)\coloneqq P_{\theta|_V}\bigl[\operatorname{Tr}_V(\varphi)\bigr]
  \in\PSH(V,\theta|_V)_{>0},
\]
as in \cite[Definition~8.1.2]{SGPT}. The two facts about the trace that we
need are the following.

The following result is recalled from \cite[Theorem~8.3.1]{SGPT}.

\begin{theorem}\label{nk:thm:trace-formula}
Let $\alpha$ be a big class, let $\theta\in\alpha$ be a smooth
representative, and let $V\subset X$ be a closed connected submanifold of
dimension $k>0$ with $V\not\subset{E_{\mathrm{nn}}}(\alpha)$. Then
\[
  \langle\alpha^k\rangle_{X\mid V}
  =\int_V\bigl(\theta|_V+\ddc\operatorname{Tr}^\theta_V(V_\theta)\bigr)^k,
\]
where the right-hand side is understood to be $0$ if
$\operatorname{Tr}^\theta_V(V_\theta)$ is not defined.
\end{theorem}

\begin{proof}
For $\epsilon>0$, set
$T_{\min,\epsilon}\coloneqq \theta+\epsilon\omega+\ddc V_{\theta+\epsilon\omega}$.
Since $V\not\subset{E_{\mathrm{nn}}}(\alpha)\supset{E_{\mathrm{nK}}}(\alpha+\epsilon\{\omega\})$
and $V_{\theta+\epsilon\omega}$ is locally bounded outside
${E_{\mathrm{nK}}}(\alpha+\epsilon\{\omega\})$, the restriction
$T_{\min,\epsilon}|_V$ is well defined. By
\cite[Remark~2.2 and equation~(2.3)]{CT22Restricted},
\[
  \langle\alpha^k\rangle_{X\mid V}
  =\lim_{\epsilon\downarrow0}\int_V(T_{\min,\epsilon}|_V)^k,
\]
the products on the right being non-pluripolar. By the equality of the
second and the last displayed quantities in
\cite[Theorem~8.3.1]{SGPT}, an identity between non-pluripolar
masses whose proof there does not involve the definition of the
restricted volume of $\alpha$, this limit equals
the right-hand side of the displayed formula, with the same convention
when the trace is not defined.
\end{proof}

The following result is recalled from \cite[Proposition~8.2.3]{SGPT}.

\begin{theorem}\label{nk:thm:trace-multiplier}
Let $V\subset X$ be a closed connected submanifold and let $\varphi$ be a
quasi-psh function on $X$ with $\nu(\varphi,V)=0$. Then, for every
$\lambda>0$,
\[
  \mathcal I\bigl(\lambda\operatorname{Tr}_V(\varphi)\bigr)\subset\NkRes{V}\mathcal I(\lambda\varphi).
\]
\end{theorem}

Since the multiplier ideals of a quasi-psh function depend only on its
$\mathcal I$-equivalence class, and since $P$-equivalent functions are
$\mathcal I$-equivalent, the inclusion of \cref{nk:thm:trace-multiplier} holds for
every trace of $\varphi$, and in particular for
$\operatorname{Tr}^\theta_V(\varphi)$ whenever the latter is defined. These results
originate in \cite{DX24}, where the trace operator was introduced.

\section{Auxiliary results}\label{nk:sec:aux}

\subsection{A bigness criterion}

We shall use the following criterion of Xiao, a weak form of the
transcendental holomorphic Morse inequalities; it is the special case of
\cite[Theorem~1.4]{Xiao15}, stated there for hermitian metrics $\omega$
with $\partial\bar\partial\omega^j=0$, in which the metric is K\"ahler.
The sharper criterion, with $4k$ replaced by $k$, was proved by
Popovici \cite[Theorem~1.1]{Pop16}, confirming a conjecture of
Boucksom--Demailly--P\u{a}un--Peternell in the K\"ahler case; the weak
form suffices for our purposes.

\begin{theorem}\label{nk:thm:nef-difference}
Let $Z$ be a connected compact K\"ahler manifold of dimension $k\geq1$ and let
$A,B$ be nef real $(1,1)$-classes on $Z$. If
$A^k-4kA^{k-1}\cdot B>0$, then $A-B$ is big.
\end{theorem}
\begin{proof}
We give the argument through Gauduchon duality and the smooth
Monge--Amp\`ere equation, following the estimate of \cite{Xiao15}.
If $k=1$, nefness gives $B\geq0$ as an intersection number, so
$\int_Z(A-B)>0$. The Poisson equation gives a positive smooth
representative of this class. We henceforth assume $k\geq2$.

Fix a K\"ahler form $\omega_Z$. By continuity of intersection numbers,
choose $\eta>0$ so small that, with $B'=B+\eta\{\omega_Z\}$,
\[
 A^k>4kA^{k-1}\cdot B'.
\]
For each sufficiently small $\varepsilon>0$, nefness gives K\"ahler
forms $a_\varepsilon\in A+\varepsilon\{\omega_Z\}$ and
$b_\varepsilon\in B'+\varepsilon\{\omega_Z\}$. Set
\[
 c_\varepsilon=(A+\varepsilon\{\omega_Z\})^k,
 \qquad
 N_\varepsilon=(A+\varepsilon\{\omega_Z\})^{k-1}
                    \cdot(B'+\varepsilon\{\omega_Z\}).
\]
Choose and fix $\varepsilon$ so small that
$c_\varepsilon>4kN_\varepsilon>0$.
We will show that $A-B'$ is pseudoeffective.

Let $\Gamma$ be any smooth strictly positive $(k-1,k-1)$-form with
$\ddc\Gamma=0$. Rescale it so that
$\int_Z b_\varepsilon\wedge\Gamma=1$, and put
$\mu=b_\varepsilon\wedge\Gamma$. This is a smooth positive probability
measure. By \cref{thm:tw-smooth}, there is a smooth function $v$ such that
\[
 \widetilde a=a_\varepsilon+\ddc v>0,
 \qquad \widetilde a^k=c_\varepsilon\mu.
\]
Indeed, the initially unknown normalization constant is
$c_\varepsilon$ by Stokes' theorem. Closedness also gives
\[
 \int_Z\widetilde a^{k-1}\wedge b_\varepsilon=N_\varepsilon.
\]
Thus the smooth positive function
\[
 H=\frac{\widetilde a^{k-1}\wedge b_\varepsilon}{\mu}
\]
has $\mu$-integral $N_\varepsilon$. The set
$E=\{H>2N_\varepsilon\}$ has $\mu(E)\leq\tfrac12$.

At a point of $Z\setminus E$, let $\lambda_1,\ldots,\lambda_k$ be
the eigenvalues of $\widetilde a$ relative to $b_\varepsilon$. The
identity for the ratio of wedge powers gives
\[
 \frac1k\sum_{i=1}^k\lambda_i^{-1}
 =\frac{\widetilde a^{k-1}\wedge b_\varepsilon}{\widetilde a^k}
 =\frac{H}{c_\varepsilon}
 \leq\frac{2N_\varepsilon}{c_\varepsilon}.
\]
Each summand is positive, hence every
$\lambda_i\geq c_\varepsilon/(2kN_\varepsilon)$. Consequently,
\[
 \widetilde a\geq
 \frac{c_\varepsilon}{2kN_\varepsilon}b_\varepsilon
 \quad\text{on }Z\setminus E.
\]
Since $\Gamma>0$ and $\ddc\Gamma=0$, integration yields
\[
 \begin{split}
 \int_Z a_\varepsilon\wedge\Gamma
 &=\int_Z\widetilde a\wedge\Gamma\\
 &\geq\frac{c_\varepsilon}{2kN_\varepsilon}\mu(Z\setminus E)
 \geq\frac{c_\varepsilon}{4kN_\varepsilon}>1
 =\int_Zb_\varepsilon\wedge\Gamma.
 \end{split}
\]
The form $a_\varepsilon-b_\varepsilon$ represents the fixed class
$A-B'$: the two $\varepsilon\{\omega_Z\}$ terms cancel. The displayed
inequality holds for every Gauduchon test form, so \cref{thm:lamari}
provides a closed positive current $S\in A-B'$. Then
$S+\eta\omega_Z\in A-B$ dominates $\eta\omega_Z$, proving bigness.
\end{proof}

The Gauduchon form serves two purposes in this proof: it tests
pseudoeffectivity, and it defines the right-hand side of the
Monge--Amp\`ere equation. Closedness makes $c_\varepsilon$ and
$N_\varepsilon$ independent of that test. This is the precise point
at which this argument needs more than a Gauduchon reference metric.

\subsection{Divisorial domination}

The following lemma converts a lower bound on the non-pluripolar mass of a
potential into a K\"ahler current whose singularities dominate those of the
potential along every prime divisor, with a positivity constant which is
independent of the potential.

\begin{lemma}\label{nk:lem:domination}
Let $(Z,\omega_Z)$ be a connected compact K\"ahler manifold of
dimension $k$, let
$\gamma$ be a smooth closed real $(1,1)$-form on $Z$ whose class is
big, and let $h_0>0$. There exists $\delta_*>0$, depending only on
$(Z,\omega_Z)$, $\gamma$ and $h_0$, with the following property. If
$0\le\eta\le1$ and
\begin{equation}\label{nk:eq:dom-hyp}
  v\in\PSH(Z,\gamma+\eta\omega_Z),\qquad
  \int_Z(\gamma+\eta\omega_Z)_v^k\ge h_0,
\end{equation}
then there is a closed positive current
$S_v\in\{\gamma\}+\eta\{\omega_Z\}$ such that
\begin{equation}\label{nk:eq:dom-concl}
  S_v\ge\delta_*\omega_Z
  \qquad\text{and}\qquad S_v\preceq_{\mathcal I} v.
\end{equation}
\end{lemma}

\begin{proof}
\emph{Step 1: analytic approximation.}
Fix $\eta$ and $v$ as in \eqref{nk:eq:dom-hyp}, and put
$\gamma_\eta\coloneqq \gamma+\eta\omega_Z$. By \cite[Theorem~1.6.2]{SGPT}
there is a quasi-equisingular approximation $(v_j)_{j\ge1}$ of $v$
consisting of potentials with analytic singularities,
\[
  v_j\in\PSH(Z,\gamma_j),\qquad
  \gamma_j\coloneqq \gamma+(\eta+\epsilon_j)\omega_Z,\qquad
  \epsilon_j\downarrow0,\qquad v_j\searrow v.
\]
Since $v\in\PSH(Z,\gamma_\eta)_{>0}$, \cite[Corollary~7.1.2]{SGPT}
gives $v_j\xrightarrow{{d_S}}P_{\gamma_\eta}[v]_\mathcal I$, so that by
\cref{nk:sec:npp}\ref{nk:it:dS-lelong} and
\cite[Propositions~3.2.1 and~3.2.6]{SGPT},
\begin{equation}\label{nk:eq:dom-lelong-conv}
  \lim_{j\to\infty}\nu(v_j,E)=\nu(v,E)
  \qquad\text{for every prime divisor $E$ over $Z$}.
\end{equation}
Since $v\le v_j$, the monotonicity theorem
(\cref{nk:sec:npp}\ref{nk:it:monotone}) and the multilinearity of non-pluripolar
products give
\begin{equation}\label{nk:eq:dom-mass}
  \int_Z(\gamma_j)_{v_j}^k
  \ge\int_Z(\gamma_j)_{v}^k
  =\int_Z\bigl((\gamma_\eta)_v+\epsilon_j\omega_Z\bigr)^k
  \ge\int_Z(\gamma_\eta)_v^k\ge h_0.
\end{equation}

\emph{Step 2: the numerical estimate.}
Let $\rho_j\colon Z_j\to Z$ be a log resolution of $v_j$, and write
the Siu decomposition of the pullback as
\[
  \rho_j^*\bigl((\gamma_j)_{v_j}\bigr)=R_j+[F_j],
\]
where $F_j$ is an effective $\mathbf Q$-divisor and $R_j$ is a closed
positive current with locally bounded potentials. Put $A_j\coloneqq \{R_j\}$.
The class $A_j$ is nef by \cref{nk:lem:residual-nef}. Since the
non-pluripolar product puts no mass on $F_j$, Bedford--Taylor theory and
\cref{nk:sec:npp}\ref{nk:it:invariance} give, together with
\eqref{nk:eq:dom-mass},
\begin{equation}\label{nk:eq:dom-Ak}
  A_j^k=\int_{Z_j}R_j^k
  =\int_{Z_j}\bigl(\rho_j^*(\gamma_j)_{v_j}\bigr)^k
  =\int_Z(\gamma_j)_{v_j}^k\ge h_0.
\end{equation}
Choose $C^\dagger>0$ so large that $\gamma+C^\dagger\omega_Z$ is
K\"ahler, and set
$K_j\coloneqq \rho_j^*\{\gamma+(\eta+\epsilon_j+C^\dagger)\omega_Z\}$. Then
$K_j$ is nef, being the pullback of a K\"ahler class, and
\[
  K_j-A_j=\{F_j\}+C^\dagger\rho_j^*\{\omega_Z\}
\]
is pseudoeffective. Since the intersection number of a pseudoeffective
class with nef classes is nonnegative, expanding
$K_j^{k-1}-A_j^{k-1}=\sum_{i=0}^{k-2}K_j^i\cdot(K_j-A_j)\cdot A_j^{k-2-i}$
and using the projection formula yields the uniform bound
\begin{equation}\label{nk:eq:dom-intersection}
  A_j^{k-1}\cdot\rho_j^*\{\omega_Z\}
  \le K_j^{k-1}\cdot\rho_j^*\{\omega_Z\}
  =\bigl(\{\gamma\}+(\eta+\epsilon_j+C^\dagger)\{\omega_Z\}\bigr)^{k-1}
   \cdot\{\omega_Z\}
  \le C_*,
\end{equation}
where
$C_*\coloneqq \max_{0\le s\le 2+C^\dagger}\bigl(\{\gamma\}+s\{\omega_Z\}\bigr)^{k-1}\cdot\{\omega_Z\}$
depends only on $(Z,\omega_Z)$ and $\gamma$; here we used
$\eta\le1$ and $\epsilon_j\le1$ for $j$ large. Now choose
$\delta_*>0$ such that
\begin{equation}\label{nk:eq:dom-delta}
  8k\delta_*C_*<h_0;
\end{equation}
this choice depends only on $(Z,\omega_Z)$, $\gamma$ and $h_0$. For
$j$ so large that $\epsilon_j<\delta_*$, the estimates
\eqref{nk:eq:dom-Ak}, \eqref{nk:eq:dom-intersection} and \eqref{nk:eq:dom-delta}
give
\[
  A_j^k-4k(\epsilon_j+\delta_*)A_j^{k-1}\cdot\rho_j^*\{\omega_Z\}
  \ge h_0-8k\delta_*C_*>0.
\]
By \cref{nk:thm:nef-difference}, applied on the compact K\"ahler manifold
$Z_j$ to the nef classes $A_j$ and
$(\epsilon_j+\delta_*)\rho_j^*\{\omega_Z\}$, the class
$A_j-(\epsilon_j+\delta_*)\rho_j^*\{\omega_Z\}$
is big. Choose a K\"ahler current $L_j$ in this class and set
\[
  \widetilde S_j\coloneqq L_j+\delta_*\rho_j^*\omega_Z+[F_j]
  \ge\delta_*\rho_j^*\omega_Z .
\]
Its cohomology class is
\[
  \{\widetilde S_j\}
  =A_j-(\epsilon_j+\delta_*)\rho_j^*\{\omega_Z\}
     +\delta_*\rho_j^*\{\omega_Z\}+\{F_j\}
  =\rho_j^*\{\gamma_j\}-\epsilon_j\rho_j^*\{\omega_Z\}
  =\rho_j^*\{\gamma_\eta\},
\]
where we used $\rho_j^*\{\gamma_j\}=A_j+\{F_j\}$.

\emph{Step 3: descent and passage to the limit.}
Since $\{\widetilde S_j\}=\rho_j^*\{\gamma_\eta\}$ and $Z_j$ is a
compact K\"ahler manifold, the $\ddc$-lemma gives
$\widetilde S_j=\rho_j^*\gamma_\eta+\ddc\widetilde s_j$ for a quasi-psh
function $\widetilde s_j$ on $Z_j$, and the lower bound
$\widetilde S_j\ge\delta_*\rho_j^*\omega_Z$ means that
$\widetilde s_j\in
\PSH\bigl(Z_j,\rho_j^*(\gamma_\eta-\delta_*\omega_Z)\bigr)$. By
\cite[Proposition~1.5.3]{SGPT}, the pullback of potentials is a
bijection from $\PSH(Z,\gamma_\eta-\delta_*\omega_Z)$ onto
$\PSH\bigl(Z_j,\rho_j^*(\gamma_\eta-\delta_*\omega_Z)\bigr)$; hence
$\widetilde s_j=s_j\circ\rho_j$ for some
$s_j\in\PSH(Z,\gamma_\eta-\delta_*\omega_Z)$. The current
$S_j\coloneqq \gamma_\eta+\ddc s_j\in\{\gamma\}+\eta\{\omega_Z\}$ then
satisfies
$S_j\ge\delta_*\omega_Z$ and $\rho_j^*S_j=\widetilde S_j$.

Since $R_j$ has locally bounded potentials and
$\widetilde S_j\ge[F_j]$, \cite[Proposition~6.1.2]{SGPT} gives
$\rho_j^*S_j=\widetilde S_j
\preceq_{\mathcal I} R_j+[F_j]=\rho_j^*\bigl((\gamma_j)_{v_j}\bigr)$, and the
bimeromorphic invariance of the $\mathcal I$-preorder
\cite[Lemma~6.1.4]{SGPT} yields $S_j\preceq_{\mathcal I} v_j$.

Normalize $\sup_Zs_j=0$. By weak compactness, a subsequence of
$(s_j)_j$ converges in $L^1(Z)$ to a quasi-psh function $s_v$, and
we set $S_v\coloneqq \gamma_\eta+\ddc s_v$. The lower bound
$S_j\ge\delta_*\omega_Z$ passes to the limit, so
$S_v\ge\delta_*\omega_Z$. For every prime divisor $E$ over $Z$,
pulling back to a model realizing $E$, upper semicontinuity of
Lelong numbers and \eqref{nk:eq:dom-lelong-conv} give
\[
  \nu(S_v,E)\ge\varlimsup_{j\to\infty}\nu(S_j,E)
  \ge\lim_{j\to\infty}\nu(v_j,E)=\nu(v,E).
\]
Another application of \cite[Proposition~6.1.2]{SGPT} gives
$S_v\preceq_{\mathcal I} v$, proving \eqref{nk:eq:dom-concl}. The constant $\delta_*$ was chosen in
\eqref{nk:eq:dom-delta} independently of $\eta\in[0,1]$ and of $v$.
\end{proof}

\subsection{Regularization with controlled positivity}

The next lemma is Demailly's regularization theorem \cite{Dem92,DPS01},
proved by the Bergman-kernel method, in a quantitative form: the loss of
positivity is $O(m^{-1})$, with a constant independent of the current
being regularized. The qualitative version, with an unspecified loss
$\epsilon_m\downarrow0$ but with the precise singularity type
$(m^{-1},\mathcal I(m\varphi))$, is \cite[Theorem~1.6.2]{SGPT}, and it
suffices for all the regularizations performed on $V$ below. The
$O(m^{-1})$ rate is needed exactly once, for the ambient current of
\cref{nk:sec:ambient}, whose positivity must be comparable in size to its
singularities; we include a complete proof. The loss of positivity comes only from
the gluing of the local Bergman kernels, and we only record the lower
bound on the singularities of $\varphi_m$, which is all that will be
used.

For $\delta>0$ and $N\ge1$, the \emph{regularized maximum}
$M_\delta\colon[-\infty,\infty)^N\to[-\infty,\infty)$ is defined by
\[
  M_\delta(t_1,\dots,t_N)\coloneqq \int_{\mathbf R^N}\max_i(t_i+\tau_i)
  \prod_{i=1}^N\frac1\delta\,\vartheta\Bigl(\frac{\tau_i}\delta\Bigr)
  \,\mathrm{d}\tau_1\cdots \mathrm{d}\tau_N,
\]
where $\vartheta\in C^\infty_c((-1,1))$ is a fixed even nonnegative
function with $\int\vartheta=1$. It has the following well-known
properties \cite[Chapter~I, Lemma~5.18]{DemCADG}: it is smooth, convex and
nondecreasing in each variable on $\mathbf R^N$;
$\max_it_i\le M_\delta(t)\le\max_it_i+N\delta$;
$M_\delta(t+c\mathbf 1)=M_\delta(t)+c$; and $M_\delta(t)$ does not
depend on $t_i$ on the set where
$t_i+\delta<\max_{j\ne i}(t_j-\delta)$, where it coincides with the
regularized maximum of the remaining variables. If $f_1,\dots,f_N$ are
smooth functions with $\zeta+\ddc f_i\ge\Gamma'$ for a smooth real
$(1,1)$-form $\Gamma'$, then
$\zeta+\ddc M_\delta(f_1,\dots,f_N)\ge\Gamma'$: by the chain rule,
$\ddc M_\delta(f_1,\dots,f_N)$ equals
$\sum_i\partial_iM_\delta\,\ddc f_i$ plus a form which is positive
semidefinite because $M_\delta$ is convex, and
$\partial_iM_\delta\ge0$ with $\sum_i\partial_iM_\delta=1$.
\begin{lemma}\label{nk:lem:regularization}
Let $(Z,\omega_Z)$ be a connected compact Hermitian manifold of
dimension $d$,
and let $\zeta$ be a smooth closed real $(1,1)$-form on $Z$. There
is a constant $c>0$, depending only on $(Z,\omega_Z)$ and not on
$\zeta$, with the following property. For every
$\varphi\in\PSH(Z,\zeta)$ and every integer $m\ge1$, there is a
quasi-psh function $\varphi_m$ on $Z$ such that
\begin{enumerate}[label=(\roman*),leftmargin=*,itemsep=.2em,topsep=.3em]
\item\label{nk:it:reg-positivity}
  $\zeta+\ddc\varphi_m\ge-\dfrac{c}{m}\,\omega_Z$;
\item\label{nk:it:reg-lower} locally on $Z$,
  $\varphi_m\ge m^{-1}\psi_{\mathcal I(m\varphi)}+O(1)$.
\end{enumerate}
\end{lemma}

The conclusion is deliberately one-sided: it does not assert analytic singularities of a prescribed type or convergence of this family. The constant in the positivity bound is uniform in the closed background; the locally bounded term in the singularity bound need not be uniform in $m$ or $\varphi$.
 
\begin{proof}
Put $\gamma_0\coloneqq -\zeta$, so that $\ddc\varphi\ge\gamma_0$.
On each coordinate ball $B_j$, write $\mathrm{d}\lambda_j$ for Euclidean
volume in its chosen coordinates. The local Hilbert spaces use these
Euclidean measures. Their comparison on fixed compact overlaps involves
only the coordinate charts, independently of $\zeta$, $\varphi$ and $m$.
 
\emph{Step 1: local Bergman weights.}
Fix once and for all, depending only on $Z$, a finite cover of $Z$ by
coordinate balls $\Omega_i$, a refinement of it by coordinate balls
$B_j^+$, $j=1,\dots,J$, and concentric balls
\[
  U_j\Subset V_j\Subset B_j\Subset B_j^+
\]
such that the $U_j$ cover $Z$. Whenever
$\overline V_j\cap\overline V_l\ne\varnothing$, the union
$B_j^+\cup B_l^+$
is contained in one of the balls, say $\Omega_{jl}$. We may write
$\gamma_0=\ddc H_{jl}$ on $\Omega_{jl}$ for a smooth function $H_{jl}$.
All coordinate metrics and coordinate volume measures are uniformly comparable on the finitely many compact overlaps. These comparison constants depend only on the fixed cover and $\omega_Z$.
 
Choose $h_j\in C^\infty(B_j^+)$ with $\ddc h_j=\gamma_0$, and set
$\psi_j\coloneqq \varphi-h_j$, which is psh on $B_j^+$. For an integer
$m\ge1$, let
\[
  \mathcal H_{j,m}\coloneqq \Bigl\{f\in\mathcal O(B_j):
     \|f\|_{j,m}^2\coloneqq \int_{B_j}|f|^2\mathrm{e}^{-m\psi_j}\,\mathrm{d}\lambda_j<\infty\Bigr\},
\]
let $(\sigma_{j,m,\ell})_\ell$ be an orthonormal basis, and define
\[
  K_{j,m}\coloneqq \sum_\ell|\sigma_{j,m,\ell}|^2,\qquad
  \psi_{j,m}\coloneqq \frac1m\log K_{j,m},\qquad
  \varphi_{j,m}\coloneqq \psi_{j,m}+h_j.
\]
The function $\psi_{j,m}$ is
psh, it is smooth outside the zero set of $K_{j,m}$, where it equals
$-\infty$, and
\begin{equation}\label{nk:eq:reg-local-positivity}
  \zeta+\ddc\varphi_{j,m}\ge0\qquad\text{on }B_j.
\end{equation}
 
\emph{Step 2: a uniform comparison on overlaps.}
Fix a pair $(j,l)$ with $\overline V_j\cap\overline V_l\ne\varnothing$.
On $B_j$ and $B_l$ we may write
\[
  h_j-H_{jl}=\Re a_j,\qquad h_l-H_{jl}=\Re a_l,
\]
with $a_j\in\mathcal O(B_j)$ and $a_l\in\mathcal O(B_l)$, since the two differences
are pluriharmonic on simply connected balls. Put $\Psi\coloneqq \varphi-H_{jl}$,
a psh function on $\Omega_{jl}$. Multiplication by $\mathrm{e}^{ma_j/2}$ is an
isometry from $\mathcal H_{j,m}$ onto the Hilbert space of holomorphic
functions on $B_j$ which are square integrable against
$\mathrm{e}^{-m\Psi}\mathrm{d}\lambda_j$, and similarly on $B_l$. Denoting by $\Psi_{j,m}$ and
$\Psi_{l,m}$ the corresponding Bergman weights, computed on $B_j$ and
on $B_l$ respectively with the common weight $\Psi$, the transformation
rule for Bergman kernels gives
\[
  \varphi_{j,m}=\Psi_{j,m}+H_{jl},\qquad
  \varphi_{l,m}=\Psi_{l,m}+H_{jl}.
\]
 
We now compare the two kernels. Choose a cut-off function
$\theta_{jl}\in C^\infty_c(B_j\cap B_l)$ with
$0\le\theta_{jl}\le1$, equal to $1$ on a
neighborhood of the compact set $\overline V_j\cap\overline V_l$, let
$r>0$ be the distance from $\overline V_j\cap\overline V_l$ to the
support of $\bar\partial\theta_{jl}$, and let $R$ be the diameter of
$B_l$, both computed in the coordinates of $B_l$. Fix
$x\in\overline V_j\cap\overline V_l$.
 
Let $f$ be in the unit ball of the $\Psi$-weighted Hilbert space on
$B_j$. Extend $\theta_{jl}f$ by zero to $B_l$ and solve
\begin{equation}\label{nk:eq:reg-dbar}
  \bar\partial g=\bar\partial(\theta_{jl}f)\qquad\text{on }B_l
\end{equation}
with respect to the weight
\[
  \Phi_{m,x}\coloneqq m\Psi+2d\log|z-x|+|z|^2 .
\]
The right-hand side of
\eqref{nk:eq:reg-dbar} is $\bar\partial$-closed, and it is supported where
$|z-x|\ge r$. Since $\Phi_{m,x}$ is psh with
$\ddc\Phi_{m,x}\ge\ddc|z|^2$, H\"ormander's $L^2$-estimate \cite[Chapter~VIII, \S6]{DemCADG}, gives a solution $g$ with
\begin{equation}\label{nk:eq:reg-hormander}
  \int_{B_l}|g|^2\mathrm{e}^{-\Phi_{m,x}}\,\mathrm{d}\lambda
  \le C_H\int_{B_l}|\bar\partial\theta_{jl}|^2|f|^2\mathrm{e}^{-\Phi_{m,x}}\,\mathrm{d}\lambda,
\end{equation}
where $\mathrm{d}\lambda$ is the Lebesgue measure and $C_H$ depends only on
$d$. Here and below, the uniform comparison between $\mathrm{d}\lambda$
and the relevant coordinate measures is absorbed into the constants. On the support of
$\bar\partial\theta_{jl}$ the factor $|z-x|^{-2d}$ is bounded by
$r^{-2d}$, so the right-hand side of \eqref{nk:eq:reg-hormander} is at most
a constant $C_1$ independent of $m$, $\Psi$ and $x$. Since
$|z-x|\le R$ on $B_l$, the left-hand side dominates the ordinary
weighted norm, and we get
\[
  \int_{B_l}|g|^2\mathrm{e}^{-m\Psi}\,\mathrm{d}\lambda_l\le C_2 .
\]
Near $x$ the function $g$ is holomorphic, and $\Psi$ is bounded
above there; the integrability of $|g|^2|z-x|^{-2d}$ near $x$ therefore
forces $g(x)=0$. Hence
$F\coloneqq \theta_{jl}f-g$ is holomorphic on $B_l$, satisfies $F(x)=f(x)$, and
has $\Psi$-weighted norm at most a uniform constant. The extremal
characterization of the Bergman kernel, applied to $F$, and the same
argument with $j$ and $l$ interchanged, yield
\begin{equation}\label{nk:eq:reg-overlap}
  \varphi_{j,m}\le\varphi_{l,m}+\frac{C_B}{m},\qquad
  \varphi_{l,m}\le\varphi_{j,m}+\frac{C_B}{m}
  \qquad\text{on }\overline V_j\cap\overline V_l,
\end{equation}
where $C_B\ge0$, the maximum of the corresponding constants over the
finitely many pairs $(j,l)$, depends only on the fixed cover
and the background geometry.
 
\emph{Step 3: global gluing.}
Choose $\chi_j\in C^\infty_c(V_j)$ with $0\le\chi_j\le1$ and
$\chi_j=1$ on $U_j$, and let $C_\chi$ be such that
$\ddc\chi_j\ge-C_\chi\omega_Z$ for all $j$. Put
$\delta\coloneqq 1/(4m)$, and set
\begin{equation}\label{nk:eq:reg-glue}
  \varphi_m(z)\coloneqq M_\delta\Bigl(\varphi_{j,m}(z)+\frac{C_B+1}m\chi_j(z)
   \ :\ z\in V_j\Bigr),
\end{equation}
the regularized maximum of the branches defined at $z$. We check that
this is a well-defined quasi-psh function on $Z$. Let
$q\in\partial V_j$ and choose $l$ with $q\in U_l$. On a small
neighborhood of $q$ we have $\chi_j=0$ and $\chi_l=1$, so that on
its intersection with $V_j\cap V_l$, \eqref{nk:eq:reg-overlap} gives
\[
  \varphi_{j,m}+\frac{C_B+1}m\chi_j+\delta
  \le\varphi_{l,m}+\frac{C_B}{m}+\delta
  <\varphi_{l,m}+\frac{C_B+1}m\chi_l-\delta
\]
off the zero locus $Z_0$ of $\mathcal I(m\varphi)$, where all branches
are finite by Step~4 below. By the independence property of
$M_\delta$, the $j$-th branch therefore does
not contribute to \eqref{nk:eq:reg-glue} near $\partial V_j\setminus Z_0$;
on $Z_0$ every branch equals $-\infty$, so \eqref{nk:eq:reg-glue}
equals $-\infty$ there whichever branches are used. In either case the local
definitions agree when branches are added or removed, and near every point
of $Z$ the function $\varphi_m$ is the regularized maximum of finitely
many branches defined on a neighborhood of that point. Each branch is
smooth where it is finite, and it satisfies
\[
  \zeta+\ddc\Bigl(\varphi_{j,m}+\frac{C_B+1}m\chi_j\Bigr)
  \ge-\frac{(C_B+1)C_\chi}{m}\omega_Z
\]
there, by \eqref{nk:eq:reg-local-positivity}. By Step~4 below, the branches
defined near a given point all equal $-\infty$ exactly on the zero
locus $Z_0$ of $\mathcal I(m\varphi)$, a proper analytic subset. Hence
$\varphi_m$ is smooth on $Z\setminus Z_0$, where it satisfies
$\zeta+\ddc\varphi_m\ge-((C_B+1)C_\chi/m)\omega_Z$ by the property of
$M_\delta$ recalled before the lemma; it is locally bounded above, and
it tends to $-\infty$ along $Z_0$, since
$M_\delta(t)\le\max_it_i+N\delta$. By the Grauert--Remmert extension
theorem for psh functions bounded above near an analytic subset
\cite{GR56} (see \cite[Theorem~1.2.1]{SGPT}), $\varphi_m$ is
quasi-psh on $Z$ and
the inequality $\zeta+\ddc\varphi_m\ge-((C_B+1)C_\chi/m)\omega_Z$ holds on
all of $Z$. Here no local potential for $\omega_Z$ is needed. On a smaller coordinate ball, add a sufficiently large quadratic function to obtain a psh function off $Z_0$, and extend it across $Z_0$. The extended current has the stated smooth lower bound off $Z_0$; any part it charges on $Z_0$ is positive, so the lower bound holds globally. This proves \ref{nk:it:reg-positivity} with
$c\coloneqq (C_B+1)C_\chi$,
which depends only on the fixed cover, the cut-off functions and
$\omega_Z$, and not on $\zeta$, $\varphi$ or $m$.
 
\emph{Step 4: the singularity ideal and the lower bound.}
Let $\mathfrak J_{j,m}$ be the ideal sheaf on $V_j$ generated by the
elements of $\mathcal H_{j,m}$. By definition of $\mathcal H_{j,m}$,
$\mathfrak J_{j,m}\subset\mathcal I(m\psi_j)=\mathcal I(m\varphi)$ on $V_j$,
the equality holding because $h_j$ is smooth. Conversely, since
$\psi_j$ is psh on $B_j^+$, and $B_j\Subset B_j^+$ is a relatively
compact Stein open subset with a basis of Stein neighborhoods,
\cite[Proposition~1.2]{DEL00}, applied with the weight
$\tfrac m2\psi_j$ and the Euclidean K\"ahler volume form $\mathrm{d}\lambda_j$ (the
multiplier ideal is defined there with the factor
$\mathrm{e}^{-2\varphi}$), shows that $\mathcal I(m\varphi)$ is generated on
$B_j$ by the orthonormal basis $(\sigma_{j,m,\ell})_\ell$ of
$\mathcal H_{j,m}$. Hence
\begin{equation}\label{nk:eq:reg-ideal-equality}
  \mathfrak J_{j,m}=\mathcal I(m\varphi)\qquad\text{on }V_j ,
\end{equation}
and in particular the zero set of $K_{j,m}$ is the zero locus of
$\mathcal I(m\varphi)$ in $V_j$, as used in Step~3.
 
We can now prove \ref{nk:it:reg-lower}. Fix $x\in Z$ and choose $j$ with
$x\in U_j$. By coherence and \eqref{nk:eq:reg-ideal-equality}, finitely
many of the $\sigma_{j,m,\ell}$, say with $\ell\le N$, generate
$\mathcal I(m\varphi)$ on a neighborhood of $x$. Therefore
$K_{j,m}\ge\sum_{\ell\le N}|\sigma_{j,m,\ell}|^2$ gives
$\varphi_{j,m}\ge m^{-1}\psi_{\mathcal I(m\varphi)}+O(1)$
on a neighborhood of $x$, the bounded term absorbing $h_j$. Since
$\chi_j\ge0$ and the regularized maximum dominates each of its
entries, $\varphi_m\ge\varphi_{j,m}$ on $V_j$, and
\ref{nk:it:reg-lower} follows.
\end{proof}

The gluing of Step~3 also produces global quasi-psh functions with a
prescribed singularity ideal.

\begin{lemma}\label{nk:lem:ideal-function}
Let $(Z,\omega_Z)$ be a connected compact Hermitian manifold and let
$\mathfrak a\subset\mathcal O_Z$ be a nonzero coherent ideal. There is a
quasi-psh function
$\Psi_{\mathfrak a}$ on $Z$ with analytic singularities of type
$(1,\mathfrak a)$ which is smooth on the complement of the zero locus of
$\mathfrak a$.
\end{lemma}

\begin{proof}
Fix a finite cover of $Z$ by coordinate balls $B_j$ on which
$\mathfrak a$ is generated by $g_{j1},\dots,g_{jN_j}$, together with
concentric balls $U_j\Subset V_j\Subset B_j$ such that the $U_j$ cover
$Z$, and put $\psi_j\coloneqq \log\sum_i|g_{ji}|^2$, a psh function on $B_j$
which is smooth outside the zero locus of $\mathfrak a$, where it equals
$-\infty$. Since the $g_{ji}$ and the $g_{li}$ generate the same
ideal, $|\psi_j-\psi_l|\le C_B$ on $\overline V_j\cap\overline V_l$ for
some constant $C_B$. Now define $\Psi_{\mathfrak a}$ by
\eqref{nk:eq:reg-glue}, with $m=1$, the branches $\psi_j$ in place of
$\varphi_{j,m}$, and the cut-off functions $\chi_j$ and the constant
$\delta=1/4$ as in Step~3 of the proof of
\cref{nk:lem:regularization}. The argument given there shows that
$\Psi_{\mathfrak a}$ is a well-defined quasi-psh function on $Z$, smooth
outside the zero locus of $\mathfrak a$, with
$\ddc\Psi_{\mathfrak a}\ge-(C_B+1)C_\chi\omega_Z$; and the bounds
$\max_it_i\le M_\delta(t)\le\max_it_i+N\delta$, together with
$0\le\chi_j\le1$ and $\psi_j=\psi_{\mathfrak a}+O(1)$, give
$\Psi_{\mathfrak a}=\psi_{\mathfrak a}+O(1)$ locally.
\end{proof}

\subsection{Relative principalization}

Let $Y$ be a complex manifold, let $W\subset Y$ be a closed connected
submanifold, and let $f\colon Y'\to Y$ be a composition of blowups with
smooth centers, each of which is either a proper analytic subset of the
current strict transform of $W$ or disjoint from it; we then call $f$
a \emph{sequence of blowups adapted to $W$}. An induction on the
number of blowups shows that the image in $Y$ of each center meets
$W$ in a proper analytic subset. Consequently $f$ is an isomorphism
over the complement of a proper analytic subset of $W$; the strict
transform $W'$ of $W$ is a connected submanifold of $Y'$ mapping
onto $W$; and no $f$-exceptional prime divisor contains $W'$,
because every such divisor maps into the image of one of the centers.

\begin{lemma}\label{nk:lem:relative-principalization}
Let $Y$ be a connected compact complex manifold, let $W\subsetneq Y$ be a
closed connected submanifold, and let $\mathfrak a\subset\mathcal O_Y$ be a coherent ideal
with $\mathfrak a\not\subset\mathcal I_W$, equivalently
$\NkRes W\mathfrak a\ne0$. There exist a sequence of blowups $f\colon Y'\to Y$
adapted to $W$ and an effective divisor $H$ on $Y'$, supported on the
exceptional locus of $f$, such that, with $W'$ the strict transform of
$W$,
\begin{enumerate}[label=\textup{(\roman*)},leftmargin=*,itemsep=.2em,topsep=.3em]
\item\label{nk:it:rp-ideal-W} $\mathcal I_W\mathcal O_{Y'}=\mathcal O_{Y'}(-H)\cdot\mathcal I_{W'}$;
\item\label{nk:it:rp-principal}
  $\mathfrak a\mathcal O_{Y'}=\mathcal O_{Y'}(-H-H')$ for an effective divisor
  $H'$ on $Y'$ which is disjoint from $W'$;
\item\label{nk:it:rp-restriction}
  $(\NkRes{W}\mathfrak a)\mathcal O_{W'}=\mathcal O_{W'}(-H|_{W'})$.
\end{enumerate}
\end{lemma}

Since no component of $H$ contains $W'$, property \ref{nk:it:rp-ideal-W}
determines $H$: it is the divisor of the ideal $\mathcal I_W\mathcal O_{Y'}$ on
$Y'\setminus W'$. If $g\colon Y''\to Y'$ is a further sequence of
blowups adapted to $W'$, with divisor $H''$ satisfying
\ref{nk:it:rp-ideal-W} for $W'$, then $g^*H+H''$ satisfies
\ref{nk:it:rp-ideal-W} for the composition $f\circ g$ and $W$.

In the proof below, we shall freely use the terminologies introduced in \cite{BM08,Wlo05}.

\begin{proof}
Put $\mathfrak b\coloneqq \NkRes W\mathfrak a\ne0$. By \cite[Theorem~1.3]{BM08}
there is a resolution of the marked ideal $(\mathfrak b,1)$ on $W$: a
finite sequence of blowups of $W$ with smooth centers, each contained
in the cosupport of the current controlled transform of $\mathfrak b$,
after which the controlled transform is the unit ideal. Blowing up the
same centers in the ambient manifold produces a sequence
$f_i\colon Y_i\to Y_{i-1}$, $i=1,\dots,t$, with $Y_0=Y$, whose
$i$-th center $C_i$ is a proper analytic subset of the strict
transform $W_{i-1}$ of $W$. (A center may have codimension one in
$W_{i-1}$; the induced map $W_i\to W_{i-1}$ is then an isomorphism,
with $E_i|_{W_i}=C_i$.) Write $E_i\subset Y_i$ for the exceptional
divisor and set
\[
  H_0=0,\quad H_i=f_i^*H_{i-1}+E_i,\qquad
  \mathfrak a_0=\mathfrak a,\quad
  \mathfrak a_i=\mathcal O_{Y_i}(E_i)\,\mathfrak a_{i-1}\mathcal O_{Y_i};
\]
here $\mathfrak a_i\subset\mathcal O_{Y_i}$ is indeed an ideal, because
$\mathfrak a_{i-1}\subset\mathcal I_{C_i}$, the restriction of
$\mathfrak a_{i-1}$ to $W_{i-1}$ vanishing on $C_i\subset W_{i-1}$.
We claim that, inductively,
\begin{equation}\label{nk:eq:rp-induction}
  \mathcal I_W\mathcal O_{Y_i}=\mathcal O_{Y_i}(-H_i)\,\mathcal I_{W_i},\qquad
  \mathfrak a\mathcal O_{Y_i}=\mathcal O_{Y_i}(-H_i)\,\mathfrak a_i,\qquad
  \NkRes{W_i}\mathfrak a_i=\mathfrak b_i,
\end{equation}
where $\mathfrak b_i$ denotes the controlled transform of
$\mathfrak b$. Indeed, the first identity follows from
$\mathcal I_{W_{i-1}}\mathcal O_{Y_i}=\mathcal O_{Y_i}(-E_i)\,\mathcal I_{W_i}$, a computation in
coordinates adapted to $C_i\subset W_{i-1}\subset Y_{i-1}$; the second
follows by pulling back its predecessor and inserting the definition of
$\mathfrak a_i$; and restricting the second to $W_i$ recovers
precisely the controlled transform
$\mathfrak b_i=\mathcal O_{W_i}(E_i|_{W_i})\,\mathfrak b_{i-1}\mathcal O_{W_i}$.

At the end, $\mathfrak b_t=\mathcal O_{W_t}$, so $\mathfrak a_t$ is the unit
ideal on a neighborhood of $W_t$. Principalize $\mathfrak a_t$ by a
second application of \cite[Theorem~1.3]{BM08}; all its centers lie in
the cosupport of $\mathfrak a_t$, hence are disjoint from the strict
transforms of $W_t$. If $H$ denotes the total transform of $H_t$
and $H'$ the divisor with $\mathfrak a_t\mathcal O_{Y'}=\mathcal O_{Y'}(-H')$, the
identities \eqref{nk:eq:rp-induction} give
\[
  \mathcal I_W\mathcal O_{Y'}=\mathcal O_{Y'}(-H)\,\mathcal I_{W'},\qquad
  \mathfrak a\mathcal O_{Y'}=\mathcal O_{Y'}(-H-H'),
\]
with $H'\cap W'=\varnothing$, proving \ref{nk:it:rp-ideal-W} and
\ref{nk:it:rp-principal}. Restricting the second identity to $W'$ gives
\ref{nk:it:rp-restriction}. All centers are adapted to $W$, and both
sequences are finite since $Y$ is compact.
\end{proof}

\Needspace{14\baselineskip}
\subsection{K\"ahler forms on blowups}

\begin{lemma}\label{nk:lem:exceptional-metric}
Let $(Z,\omega_Z)$ be a compact K\"ahler manifold and let
$\rho\colon Z'\to Z$ be a composition of blowups with smooth centers.
There are an effective $\rho$-exceptional $\mathbf Q$-divisor $E_\rho$
on $Z'$ and a smooth Hermitian metric on $\mathcal O_{Z'}(E_\rho)$ such that
\[
  \rho^*\omega_Z-\Theta_{E_\rho}\quad\text{is a K\"ahler form on }Z'.
\]
If $\rho$ is an isomorphism, one may take $E_\rho=0$.
\end{lemma}

\begin{proof}
Blowups along centers of codimension one are isomorphisms and may be
omitted. For a single blowup $f\colon Z_1\to Z$ with smooth center
$C$ of codimension at least two and exceptional divisor $F$, we use
the construction of
\cite[Chapter~VII, Proposition~12.4 and its proof]{DemCADG}: one has
$F=\mathbf P(N_{C/Z})$ and
$\mathcal O_{Z_1}(-F)|_F=\mathcal O_{\mathbf P(N_{C/Z})}(1)$, so a Hermitian metric on
$N_{C/Z}$ induces a metric on $\mathcal O_{Z_1}(-F)|_F$ whose curvature is
positive along every fiber of $F\to C$; extend it smoothly to
$\mathcal O_{Z_1}(-F)$. The resulting form $-\Theta_F$ is then positive on
every nonzero vector of $\ker \mathrm{d}f$, these being exactly the vectors
tangent to the fibers of $F\to C$, while $f^*\omega_Z$ is
semipositive with kernel $\ker \mathrm{d}f$. Since $-\Theta_F$ is positive on
a neighborhood of the unit vectors of $\ker \mathrm{d}f$ and $f^*\omega_Z$ is
bounded below on the complementary compact subset of the unit tangent
bundle, $f^*\omega_Z-\epsilon\Theta_F$ is a K\"ahler form for every
sufficiently small $\epsilon>0$; note that no integrality of the class
$\{\omega_Z\}$ is required. The general case follows by induction on
the number of blowups: if
$\omega_{i-1}\coloneqq \rho_{i-1}^*\omega_Z-\Theta_{E_{i-1}}$ is a K\"ahler
form on $Z_{i-1}$ and $f_i\colon Z_i\to Z_{i-1}$ is the next blowup,
with exceptional divisor $F_i$, then
$f_i^*\omega_{i-1}-\epsilon_i\Theta_{F_i}$ is K\"ahler for a small
rational $\epsilon_i>0$, and this form equals
$\rho_i^*\omega_Z-\Theta_{E_i}$ with
$E_i\coloneqq f_i^*E_{i-1}+\epsilon_iF_i$, endowed with the induced metric.
\end{proof}

\subsection{Gluing}

The following proposition isolates the maximum argument used in the final step. The pointwise comparison on a compact collar is its substantive hypothesis.

\begin{proposition}\label{nk:prop:gluing}
Let $Y$ be a connected compact K\"ahler manifold, let $\theta_Y$ be
a smooth closed real $(1,1)$-form on $Y$, let $U_1\Subset U_2$ be
open subsets of $Y$, and let $y\in U_1$. Suppose that there are a
quasi-psh function $\varphi_G$ on $Y$ and a quasi-psh function
$\varphi_U$ on a neighborhood of $\overline U_2$ such that
\begin{enumerate}[label=\textup{(\roman*)},leftmargin=*,itemsep=.25em]
\item\label{nk:it:gl-current}
  $\theta_Y+\ddc\varphi_G$ is a K\"ahler current, and
  $\varphi_G(y)=-\infty$;
\item\label{nk:it:gl-local}
  $\theta_Y+\ddc\varphi_U$ dominates a K\"ahler form on a neighborhood
  of $\overline U_2$, and $\varphi_U$ is smooth near $y$;
\item\label{nk:it:gl-collar}
  $\varphi_U\le\varphi_G+C$ on $\overline U_2\setminus U_1$ for some
  constant $C$.
\end{enumerate}
Then there is a K\"ahler current $\mathcal T_Y\in\{\theta_Y\}$ which is
smooth near $y$. In particular, $y\notin{E_{\mathrm{nK}}}(\{\theta_Y\})$.
\end{proposition}

\begin{proof}
The function
\[
  \Phi\coloneqq 
  \begin{cases}
    \max\{\varphi_U-C,\varphi_G\}&\text{on }U_2,\\
    \varphi_G&\text{on }Y\setminus\overline U_1
  \end{cases}
\]
is well defined: on $U_2\setminus\overline U_1$, hypothesis
\ref{nk:it:gl-collar} gives $\varphi_U-C\le\varphi_G$, so the two
branches agree. By \ref{nk:it:gl-current}, \ref{nk:it:gl-local} and
compactness, there is a K\"ahler form $\omega_0$ on $Y$ with
$\theta_Y+\ddc\varphi_G\ge\omega_0$ on $Y$ and
$\theta_Y+\ddc\varphi_U\ge\omega_0$ on a neighborhood of
$\overline U_2$. The maximum of $(\theta_Y-\omega_0)$-psh functions
is again $(\theta_Y-\omega_0)$-psh, so
$\mathcal T_Y\coloneqq \theta_Y+\ddc\Phi\ge\omega_0$ is a K\"ahler current on
$Y$. Since $\varphi_U$ is smooth near $y$, whereas upper
semicontinuity and $\varphi_G(y)=-\infty$ imply that $\varphi_G$ is
arbitrarily negative near $y$, we have $\Phi=\varphi_U-C$ near
$y$. Thus $\mathcal T_Y$ is smooth there, so that
$y\notin{E_{\mathrm{nK}}}(\{\theta_Y\})$; note that the class $\{\theta_Y\}$ is
big, as it contains a K\"ahler current.
\end{proof}

The absence of isolated non-K\"ahler points was proved in
\cref{prop:bc-no-isolated-locus}. It will be used only in the last step,
when componentwise vanishing is converted into equality of the two loci.

\section{Reduction to a positive trace mass}\label{nk:sec:setup}

In this section we begin the proof of \cref{nk:thm:main}. We first reduce to
a situation in which $\alpha$ is big and $V$ is smooth, and then
express the numerical restricted volume as the mass of a trace.

\subsection{Reductions}

Let $X$, $\alpha$ and $V$ be as in \cref{nk:thm:main}, and put
$k\coloneqq \dim V>0$. We have to prove that
$\langle\alpha^k\rangle_{X\mid V}=0$.
If $\alpha$ is not big, then
${E_{\mathrm{nK}}}(\alpha)=X$, so $V=X$ and
$\langle\alpha^n\rangle_{X\mid X}=\vol(\alpha)=0$ by
\cref{nk:sec:restricted-volumes}. We may therefore assume that $\alpha$ is
big; then ${E_{\mathrm{nK}}}(\alpha)$ is a proper analytic subset of $X$, so that
$k<n$. If $V\subset{E_{\mathrm{nn}}}(\alpha)$, the desired vanishing holds by the
definition of the numerical restricted volume. We may therefore assume
that $V\not\subset{E_{\mathrm{nn}}}(\alpha)$, and we argue by contradiction: assume
that
$h\coloneqq \langle\alpha^k\rangle_{X\mid V}>0$.

\begin{lemma}\label{nk:lem:resolution}
In the situation above, there is a modification
$\mu\colon\widetilde X\to X$ which is an isomorphism over
$X\setminus V_{\mathrm{sing}}$ and such that the strict transform
$\widetilde V$ of $V$ is smooth. The triple
$(\widetilde X,\mu^*\alpha,\widetilde V)$ satisfies all the standing
assumptions: $\widetilde X$ is a connected compact K\"ahler manifold,
$\mu^*\alpha$ is big, $\widetilde V$ is a positive-dimensional
irreducible component of ${E_{\mathrm{nK}}}(\mu^*\alpha)$ not contained in
${E_{\mathrm{nn}}}(\mu^*\alpha)$, and
$\langle(\mu^*\alpha)^k\rangle_{\widetilde X\mid\widetilde V}=h>0$.
\end{lemma}

\begin{proof}
Let $\mu$ be an embedded resolution of $V\subset X$, that is, a
composition of blowups with smooth centers lying over
$V_{\mathrm{sing}}$ whose composite is an isomorphism over
$X\setminus V_{\mathrm{sing}}$ and such that the strict transform of
$V$ is smooth; it exists by Hironaka's theorem, see for instance
\cite{BM97}. Blowups of compact K\"ahler
manifolds along smooth centers are compact K\"ahler, and the pullback of
a big class under a modification is big. Since $\mu$ is an isomorphism
over $X\setminus V_{\mathrm{sing}}$, we have
$V\not\subset\mu(\operatorname{Exc}(\mu))$, so
\cref{nk:sec:restricted-volumes}\ref{nk:it:rv-birational} gives
$\langle(\mu^*\alpha)^k\rangle_{\widetilde X\mid\widetilde V}=h$. By
\eqref{nk:eq:transformation-rules},
${E_{\mathrm{nn}}}(\mu^*\alpha)=\mu^{-1}({E_{\mathrm{nn}}}(\alpha))$, so
$\widetilde V\subset{E_{\mathrm{nn}}}(\mu^*\alpha)$ would give
$V=\mu(\widetilde V)\subset{E_{\mathrm{nn}}}(\alpha)$, which is excluded.

It remains to see that $\widetilde V$ is an irreducible component of
${E_{\mathrm{nK}}}(\mu^*\alpha)=\mu^{-1}({E_{\mathrm{nK}}}(\alpha))\cup\operatorname{Exc}(\mu)$. Clearly
$\widetilde V\subset\mu^{-1}(V)\subset{E_{\mathrm{nK}}}(\mu^*\alpha)$. Let
$Z\supset\widetilde V$ be an irreducible component of
${E_{\mathrm{nK}}}(\mu^*\alpha)$. Then $\mu(Z)$ is an irreducible analytic subset of
$X$ containing $V$ and contained in
${E_{\mathrm{nK}}}(\alpha)\cup\mu(\operatorname{Exc}(\mu))\subset{E_{\mathrm{nK}}}(\alpha)\cup V_{\mathrm{sing}}
={E_{\mathrm{nK}}}(\alpha)$. Since $V$ is an irreducible component of
${E_{\mathrm{nK}}}(\alpha)$, we get $\mu(Z)=V$, hence
$Z\subset\mu^{-1}(V)=\widetilde V\cup\mu^{-1}(V_{\mathrm{sing}})$. As
$Z$ is irreducible and $\mu(Z)=V\not\subset V_{\mathrm{sing}}$, we
conclude that $Z\subset\widetilde V$, that is, $Z=\widetilde V$.
\end{proof}

Replacing $(X,\alpha,V)$ by $(\widetilde X,\mu^*\alpha,\widetilde V)$,
we are reduced to the following situation, which remains in force for the
rest of the proof.

\paragraph{Standing assumptions.}
$X$ is a connected compact K\"ahler manifold of dimension $n$ with
K\"ahler form $\omega$; $\alpha$ is a big class with a fixed smooth
representative $\theta\in\alpha$; $V\subset X$ is a connected
submanifold of dimension $k$, $0<k<n$, which is an irreducible
component of ${E_{\mathrm{nK}}}(\alpha)$ and is not contained in ${E_{\mathrm{nn}}}(\alpha)$;
and $h=\langle\alpha^k\rangle_{X\mid V}>0$. We write
$\omega_V\coloneqq \omega|_V$ and $\gamma\coloneqq \theta|_V$.


We shall reach a contradiction in \cref{nk:sec:gluing}. All modifications
constructed below are sequences of blowups adapted to $V$ in the sense
of \cref{nk:sec:aux}, hence isomorphisms over a nonempty relatively open
subset of $V$.

\subsection{The trace of \texorpdfstring{$V_\theta$}{V\_theta}}

Since $V\not\subset{E_{\mathrm{nn}}}(\alpha)$, choose
$x_*\in V\setminus{E_{\mathrm{nn}}}(\alpha)$. Then
$\nu(V_\theta,x_*)=0$ by \cref{nk:sec:nK}.
By \cref{nk:thm:trace-formula},
\[
  h=\langle\alpha^k\rangle_{X\mid V}
   =\int_V\bigl(\gamma+\ddc\operatorname{Tr}^\theta_V(V_\theta)\bigr)^k,
\]
where the right-hand side would be $0$ if the trace
$\operatorname{Tr}^\theta_V(V_\theta)$ were not defined. Since $h>0$, the trace is
defined, and we set
$u\coloneqq \operatorname{Tr}^\theta_V(V_\theta)\in\PSH(V,\gamma)_{>0}$, so that
$\int_V\gamma_u^k=h$.
Since $\PSH(V,\gamma)_{>0}\ne\varnothing$, the class
$\{\gamma\}=\alpha|_V$ is big by \cref{nk:sec:npp}\ref{nk:it:bigness}; this
is the bigness hypothesis needed in \cref{nk:lem:domination} for
$(Z,\omega_Z,\gamma)=(V,\omega_V,\gamma)$.

For every positive integer $q$, set
$\mathfrak B_q\coloneqq \mathcal I(qV_\theta)$ and denote their restrictions
$\overline{\mathfrak B}_q\coloneqq \NkRes{V}\mathfrak B_q=\operatorname{im}\bigl(\mathfrak B_q\to\mathcal O_V\bigr)$
to $V$.
By Skoda's theorem and $\nu(V_\theta,x_*)=0$, $\mathfrak B_q$ and
$\overline{\mathfrak B}_q$ are unit ideals near $x_*$, and in particular
$\overline{\mathfrak B}_q\ne0$. By \cref{nk:thm:trace-multiplier}, we have $\mathcal I(qu)\subset\overline{\mathfrak B}_q$.
We call the $\overline{\mathfrak B}_q$ the \emph{restricted ambient ideals}.

\Needspace{12\baselineskip}
\subsection{Analytic approximations of the trace}\label{nk:sec:trace-approx}

Let $(\Phi^D_j)_{j\ge1}$ be a quasi-equisingular approximation of
$V_\theta$ as in \cref{nk:sec:npp}, with $\Phi^D_j$ of analytic
singularities of type $(2^{-j},\mathcal I(2^jV_\theta))$. We index it
by the dyadic level $q=2^j$ and write $\Phi^D_q$; thus
\[
  \Phi^D_q\in\PSH(X,\theta+\eta_q\omega),\qquad
  \eta_q\downarrow0.
\]
The function $\Phi^D_q$ has analytic singularities of type
$(q^{-1},\mathfrak B_q)$ and decreases to $V_\theta$ as $q\to\infty$ along the dyadic
integers. Since $\mathfrak B_q\not\subset\mathcal I_V$, the restriction
$\phi_q\coloneqq \Phi^D_q|_V$ is not identically $-\infty$, and it has analytic
singularities of type $(q^{-1},\overline{\mathfrak B}_q)$.
Set $\beta_q\coloneqq \gamma+\eta_q\omega_V$, so that
$\phi_q\in\PSH(V,\beta_q)$.

\begin{lemma}\label{nk:lem:trace-approximants}
With the notation above, the following hold.
\begin{enumerate}[label=\textup{(\roman*)},leftmargin=*,itemsep=.2em,topsep=.3em]
\item\label{nk:it:vq-mass} $\int_V(\beta_q)_{\phi_q}^k\to h$ as
  $q\to\infty$ along the dyadic integers.
\item\label{nk:it:vq-order} $u\preceq_P\phi_q$ for every sufficiently large
  dyadic $q$.
\item\label{nk:it:vq-lower} for every integer $i\ge1$ and every local
  section $g$ of $\overline{\mathfrak B}_q^{\,i}$, the function
  $|g|^2\mathrm{e}^{-qi\,\phi_q}$ is locally bounded on $V$ outside a pluripolar set.
\end{enumerate}
\end{lemma}

\begin{proof}
Choose $A$ sufficiently large that $A>\sup_q\eta_q$ and
$\beta\coloneqq \gamma+A\omega_V$ is a K\"ahler form on $V$, so that
$\phi_q\in\PSH(V,\beta_q)\subset\PSH(V,\beta)$
for all $q$. By definition, $\phi_q\xrightarrow{d_{S}}u$.
Multilinearity \cite[Proposition~2.2.1]{SGPT} and the mixed-mass
continuity in \cref{nk:sec:npp}\ref{nk:it:dS-mass}, together with
$\beta_q\to\gamma$, now give
$\int_V(\beta_q)_{\phi_q}^k\to\int_V\gamma_u^k=h$, which is
\ref{nk:it:vq-mass}.

By \ref{nk:it:vq-mass}, multilinearity, and positivity of mixed
non-pluripolar products,
$\int_V\beta_{\phi_q}^k\ge\int_V(\beta_q)_{\phi_q}^k>0$ for every
sufficiently large dyadic $q$; also
$\int_V\beta_u^k\ge\int_V\gamma_u^k=h>0$, so
$u,\phi_q\in\PSH(V,\beta)_{>0}$. Fix such a $q$. Since
$(\Phi^D_{q'})_{q'}$ decreases, $\phi_{q'}\le\phi_q$ for all dyadic
$q'\ge q$, hence $\phi_{q'}\preceq_P\phi_q$. As
$\phi_{q'}\xrightarrow{d_{S}}u$,
\cite[Proposition~6.2.5]{SGPT}, applied in $\PSH(V,\beta)$ with the
constant net $\phi_q$, gives $u\preceq_P\phi_q$, proving
\ref{nk:it:vq-order}.

Finally, if $g$ is
a local section of $\overline{\mathfrak B}_q^{\,i}$, then $g$ is a finite sum of
products of $i$ restrictions to $V$ of local generators of $\mathfrak B_q$,
so $|g|^2\le c_g\bigl(\sum_a|f_a|_V|^2\bigr)^i=c_g'\mathrm{e}^{qi\phi_q}$ locally,
where $f_a$ are local generators of $\mathfrak B_q$ and
$\phi_q=q^{-1}\log\sum_a|f_a|_V|^2+O(1)$. This proves \ref{nk:it:vq-lower}.
\end{proof}

\section{A positive current on the submanifold}\label{nk:sec:positive-current}

We keep the assumptions and notation of \cref{nk:sec:setup}. In this section
we produce a K\"ahler current with analytic singularities on $V$ whose
multiplier ideal is compared with the restricted ambient ideals
$\overline{\mathfrak B}_m$ and $\overline{\mathfrak B}_M$ at two levels $m\ll M$.

\Needspace{15\baselineskip}
\subsection{Fixed data}\label{nk:sec:fixed-data}

Let $\delta_*>0$ be the constant provided by \cref{nk:lem:domination}
for $(Z,\omega_Z,\gamma,h_0)=(V,\omega_V,\gamma,2^{-k-1}h)$, and put
$\delta_0\coloneqq \delta_*/4$.

By \cref{nk:sec:nK}, there is a K\"ahler current with analytic singularities
\begin{equation}\label{nk:eq:def-R}
  R=\theta+\ddc\Phi_R\in\alpha,\qquad R\ge\delta_R\omega,\qquad
  \text{with polar set }{E_{\mathrm{nK}}}(\alpha),
\end{equation}
for some $\delta_R>0$. Let $Z'$ be the union of the irreducible
components of ${E_{\mathrm{nK}}}(\alpha)$ other than $V$, and put
$\mathfrak j\coloneqq \mathcal I_{Z'}\subset\mathcal O_X$,
with $\mathfrak j=\mathcal O_X$ if $Z'=\varnothing$. Since $V$ is an
irreducible component of ${E_{\mathrm{nK}}}(\alpha)$, it is not contained in $Z'$;
hence $\NkRes V\mathfrak j\ne0$.

The potential $\Phi_R$ has polar set
${E_{\mathrm{nK}}}(\alpha)=V\cup Z'=Z(\mathcal I_V\cap\mathfrak j)$. R\"uckert's
Nullstellensatz \cite[Chapter~II, Theorem~4.22]{DemCADG} and compactness therefore give
a constant $A_R>0$ such that, locally on $X$,
\begin{equation}\label{nk:eq:PhiR-lower}
  \Phi_R\ge A_R\bigl(\psi_{\mathcal I_V}+\psi_{\mathfrak j}\bigr)+O(1).
\end{equation}
Indeed, if $\Phi_R=c\psi_{\mathfrak r}+O(1)$, compactness supplies
$N$ with $\mathcal I_V^N\mathfrak j^N\subset\mathfrak r$, and one may
take $A_R=cN$.

Let $c_0$ be the constant provided by \cref{nk:lem:regularization} for
$(Z,\omega_Z)=(X,\omega)$; it depends only on $(X,\omega)$. Now
choose constants in the following order:
\begin{equation}\label{nk:eq:constants}
  \begin{alignedat}{2}
  &L\in\mathbf N &&\quad\text{with }L\delta_R>c_0,\\
  &C\in\mathbf N &&\quad\text{with }C-1\ge LA_R.
  \end{alignedat}
\end{equation}
Both constants depend only on the data fixed so far, and not on the
levels $m$ and $M$ chosen below.


For clarity, the order of choices is as follows. Each row may depend on
all preceding rows, but never on a later approximation level.
\begin{center}
\small
\begin{tabular}{p{0.22\textwidth}p{0.67\textwidth}}
\toprule
Data & What is fixed before proceeding \\
\midrule
$\delta_*,\delta_0,R$ & The trace mass and the uniform divisorial domination bound. \\
$L,C$ & The ambient regularization loss and the singularity bound for $R$. \\
$m,K,Y_{\mathrm{fix}}$ & A positive rooftop, a strict current on $V$ and a principalization of $\mathfrak B_m$. \\
$\Psi_m,c_\Psi$ & The fixed correction potential and its curvature bound. \\
$\ell,M=C\ell$ & A sufficiently fine dyadic level on $V$ and the ambient level. \\
$Y,E_{\mathrm{rel}},\tau$ & Moving principalizations and a small exceptional correction. \\
\bottomrule
\end{tabular}
\end{center}

\subsection{A rooftop envelope of positive mass}

\begin{lemma}\label{nk:lem:positive-rooftop}
For every sufficiently large dyadic $m$, there is a function
$h_C\in\PSH(V,\beta_m)$ such that
\[
  h_C\le Cu-(C-1)\phi_m\ \text{quasi-everywhere}
  \qquad\text{and}\qquad
  \int_V(\beta_m)_{h_C}^k\ge2^{-k-1}h .
\]
\end{lemma}

\begin{proof}
Put $a_m\coloneqq \int_V(\beta_m)_u^k$ and $b_m\coloneqq \int_V(\beta_m)_{\phi_m}^k$.
Then $b_m\to h$ by
\cref{nk:lem:trace-approximants}\ref{nk:it:vq-mass}, and since
$u\in\PSH(V,\gamma)$, multilinearity gives
\[
  a_m=\sum_{j=0}^k\binom kj\eta_m^{k-j}\int_V\gamma_u^j\wedge\omega_V^{k-j}
  \longrightarrow\int_V\gamma_u^k=h.
\]
In particular $(2C)^ka_m>((2C)^k-1)b_m$ for all large $m$. Moreover
$u\preceq_P\phi_m$ for $m$ large
(\cref{nk:lem:trace-approximants}\ref{nk:it:vq-order}), and
$u,\phi_m\in\PSH(V,\beta_m)_{>0}$ for $m$ large.
\Cref{nk:thm:rooftop}, applied on $V$ with $\theta=\beta_m$,
$\varphi=u$, $\psi=\phi_m$ and $t=2C$, shows that
$w\coloneqq P_{\beta_m}\bigl(2Cu-(2C-1)\phi_m\bigr)\in\PSH(V,\beta_m)_{>0}$. Set
$h_C\coloneqq \tfrac12(w+\phi_m)\in\PSH(V,\beta_m)$. Since
$w\le 2Cu-(2C-1)\phi_m$ quasi-everywhere,
$h_C\le Cu-(C-1)\phi_m$ quasi-everywhere. By multilinearity and the
positivity of mixed non-pluripolar products,
\[
  \int_V(\beta_m)_{h_C}^k
  =\int_V\Bigl(\tfrac12(\beta_m)_{w}+\tfrac12(\beta_m)_{\phi_m}\Bigr)^k
  \ge2^{-k}\int_V(\beta_m)_{\phi_m}^k=2^{-k}b_m,
\]
and $2^{-k}b_m>2^{-k-1}h$ for $m$ large, since $b_m\to h$.
\end{proof}

From now on we fix a dyadic integer $m$ so large that
\cref{nk:lem:positive-rooftop} applies, that $\eta_m\le1$, and that
$\eta_m<\delta_*/2$.
Apply \cref{nk:lem:domination} with
$(Z,\omega_Z,\gamma,h_0)=(V,\omega_V,\gamma,2^{-k-1}h)$,
$\eta=\eta_m$ and
$v=h_C$; the hypothesis \eqref{nk:eq:dom-hyp} is the mass bound of
\cref{nk:lem:positive-rooftop}. We
obtain a closed positive current $K^+\in\{\beta_m\}$ with
$K^+\ge\delta_*\omega_V$ and $K^+\preceq_{\mathcal I} h_C$. 
Since $\eta_m<\delta_*/2$, we obtain a K\"ahler
current
\begin{equation}\label{nk:eq:def-K}
  K\coloneqq K^+-\eta_m\omega_V\in\alpha|_V,\qquad
  K\ge\frac{\delta_*}2\omega_V,\qquad
  K\preceq_{\mathcal I} h_C .
\end{equation}
We write $K=\gamma+\ddc\varphi_K$ with a quasi-psh function
$\varphi_K$ on $V$.

\subsection{The fixed principalization and the auxiliary function}\label{nk:sec:fixed-principalization}

The restriction $\overline{\mathfrak B}_m$ is nonzero by
\cref{nk:sec:setup}. Apply \cref{nk:lem:relative-principalization} to
$\mathfrak B_m$ on $X$, with $V$ as submanifold. We obtain a sequence
of blowups $\pi_{{\mathrm{fix}}}\colon Y_{{\mathrm{fix}}}\to X$ adapted to $V$ and an
effective $\pi_{{\mathrm{fix}}}$-exceptional divisor $H^{{\mathrm{fix}}}_m$ on
$Y_{{\mathrm{fix}}}$ with the properties \ref{nk:it:rp-ideal-W}--\ref{nk:it:rp-restriction}
of that lemma; in particular $\mathfrak B_m$ is invertible on $Y_{{\mathrm{fix}}}$,
and we write
\begin{equation}\label{nk:eq:def-Gm-tilde}
  \mathfrak B_m\mathcal O_{Y_{{\mathrm{fix}}}}=\mathcal O_{Y_{{\mathrm{fix}}}}(-\widetilde G_m),
  \qquad \widetilde G_m\ge H^{{\mathrm{fix}}}_m .
\end{equation}
Let $W_m\subset Y_{{\mathrm{fix}}}$ be the strict transform of $V$, so that
$\pi_{{\mathrm{fix}}}$ restricts to a composition of blowups $W_m\to V$ along
smooth centers, and let $G_m\coloneqq \widetilde G_m|_{W_m}$, the
divisor of the invertible ideal $\overline{\mathfrak B}_m\mathcal O_{W_m}$.

Apply \cref{nk:lem:ideal-function} to the product ideal
$\mathfrak j\mathfrak B_m$ on $X$: we obtain a quasi-psh function
$\Psi_m$ on $X$, smooth outside the proper analytic subset
$Z'\cup Z(\mathfrak B_m)$, with
\begin{equation}\label{nk:eq:def-Psi}
  \Psi_m=\psi_{\mathfrak j}+\psi_{\mathfrak B_m}+O(1)\ \text{locally},
  \qquad
  \ddc\Psi_m\ge-c_\Psi\,\omega
\end{equation}
for some constant $c_\Psi>0$; we normalize $\Psi_m\le0$. All the data
introduced in this subsection depend on $m$ but not on the second level
$M$.

\subsection{The second level and the current \texorpdfstring{$K_\ell$}{K\_l}}

By \eqref{nk:eq:def-K}, $\varphi_K\in\PSH(V,\gamma-(\delta_*/2)\omega_V)$.
By \cite[Theorem~1.6.2]{SGPT}, $\varphi_K$ admits a
quasi-equisingular approximation consisting of potentials with analytic
singularities of type $(2^{-j},\mathcal I(2^j\varphi_K))$; as in
\cref{nk:sec:trace-approx}, we index it by the dyadic level $q=2^j$ and
write it as $(\varphi_{K,q})_q$. Its analytic singularities have type
$(q^{-1},\mathcal I(q\varphi_K))$, and
\[
  \gamma+\ddc \varphi_{K,q}\ge\Bigl(\frac{\delta_*}2-\eta'_q\Bigr)\omega_V,
\]
with $\eta'_q\downarrow 0$. We fix
this approximation once and for all; it depends only on data fixed
before the choice of $M$ below.

Now choose a dyadic integer $\ell\ge m$ so large that
$\eta'_\ell\le\delta_0$ and that $M\coloneqq C\ell$ satisfies $M>L$ and
$M\ge2(C-1)c_\Psi/\delta_0$; these conditions involve only data fixed
before this choice. Since $m$ and $\ell$ are dyadic with
$\ell\ge m$, $m$ divides $\ell$, so both $M/m$ and
$p\coloneqq (C-1)\ell/m$ are integers, and $mp=\ell(C-1)$. The
subadditivity \eqref{nk:eq:subadditivity} gives
$\mathfrak B_M\subset\mathfrak B_m^{M/m}$, hence
$\overline{\mathfrak B}_M\subset\overline{\mathfrak B}_m^{\,M/m}$. Set
\begin{equation}\label{nk:eq:def-K-ell}
  K_\ell\coloneqq \gamma+\ddc\varphi_{K,\ell}\in\alpha|_V,\qquad
  K_\ell\ge\Bigl(\frac{\delta_*}2-\eta'_\ell\Bigr)\omega_V
  \ge\delta_0\,\omega_V,
\end{equation}
the last inequality holding because $\eta'_\ell\le\delta_0$. The
potential $\varphi_{K,\ell}$ has analytic singularities of type
$(\ell^{-1},\mathfrak B)$, where $\mathfrak B\coloneqq \mathcal I(\ell\varphi_K)\subset\mathcal O_V$ and
$\ell^{-1}=C/M$.

\begin{lemma}\label{nk:lem:ideal-comparison}
The ideals $\mathfrak B$, $\overline{\mathfrak B}_m$ and $\overline{\mathfrak B}_M$ on $V$ satisfy
\begin{equation}\label{nk:eq:ideal-comparison}
  \mathfrak B\,\overline{\mathfrak B}_m^{\,p}\subset\mathcal I(Mu)\subset\overline{\mathfrak B}_M .
\end{equation}
Consequently, on any modification of $V$ on which $\overline{\mathfrak B}_m$,
$\overline{\mathfrak B}_M$ and $\mathfrak B$ are principal, with divisors $G_m$, $G_M$
and $G_B$ respectively, the divisors
\[
  G_M-\frac MmG_m\qquad\text{and}\qquad F_B\coloneqq G_B+pG_m-G_M
\]
are effective, and
\begin{equation}\label{nk:eq:GB-decomposition}
  \frac1\ell G_B=\frac1MG_M+\frac{C-1}M\Bigl(G_M-\frac MmG_m\Bigr)
     +\frac CMF_B .
\end{equation}
\end{lemma}

\begin{proof}
The second inclusion in \eqref{nk:eq:ideal-comparison} is the trace
inclusion $\mathcal I(Mu)\subset\overline{\mathfrak B}_M$ of \cref{nk:sec:setup}. For the
first one, let $f$ be
a local section of $\mathfrak B$ and $g$ a local section of
$\overline{\mathfrak B}_m^{\,p}$. By \eqref{nk:eq:def-K}, $f\in\mathcal I(\ell h_C)$, that is,
$|f|^2\mathrm{e}^{-\ell h_C}$ is locally integrable. By
\cref{nk:lem:trace-approximants}\ref{nk:it:vq-lower},
together with $mp=\ell(C-1)$, the function
$|g|^2\mathrm{e}^{-\ell(C-1)\phi_m}=|g|^2\mathrm{e}^{-mp\,\phi_m}$
is locally bounded.
By \cref{nk:lem:positive-rooftop}, $h_C\le Cu-(C-1)\phi_m$ quasi-everywhere,
hence almost everywhere, and $\ell C=M$; therefore
\[
  |fg|^2\mathrm{e}^{-Mu}\le|f|^2\mathrm{e}^{-\ell h_C}\,|g|^2\mathrm{e}^{-\ell(C-1)\phi_m}
  \qquad\text{almost everywhere},
\]
so that $fg\in\mathcal I(Mu)$. This proves \eqref{nk:eq:ideal-comparison}.

Since $\overline{\mathfrak B}_M\subset\overline{\mathfrak B}_m^{\,M/m}$, on a modification where
both ideals are principal we have $G_M\ge(M/m)G_m$. Similarly, \eqref{nk:eq:ideal-comparison} gives
$\mathfrak B\,\overline{\mathfrak B}_m^{\,p}\subset\overline{\mathfrak B}_M$, that is,
$G_B+pG_m\ge G_M$, so $F_B\ge0$. Finally, substituting
$G_B=G_M-pG_m+F_B$, $M=C\ell$ and $p=(C-1)\ell/m$ into the
right-hand side of \eqref{nk:eq:GB-decomposition}, the coefficient of
$G_M$ is $(1+(C-1)-C)/M=0$, the coefficient of $G_m$ is
$-(C-1)/m+Cp/M=0$, and the coefficient of $G_B$ is $C/M=1/\ell$.
\end{proof}

\section{Construction of the gluing data}\label{nk:sec:principalization}

We keep the notation of \cref{nk:sec:setup,nk:sec:positive-current}. In this
section we construct the data to which \cref{nk:prop:gluing} will be
applied in \cref{nk:sec:gluing}.

\subsection{The ambient current}\label{nk:sec:ambient}

Apply \cref{nk:lem:regularization} to $V_\theta$ on $(X,\omega)$ with
$\zeta=\theta$ and the level $M$: after an additive
normalization we obtain a quasi-psh function $U_M\le0$ on $X$ with
\begin{equation}\label{nk:eq:def-TM}
  T_M\coloneqq \theta+\ddc U_M\in\alpha,\qquad T_M\ge-\frac{c_0}M\omega,\qquad
  U_M\ge\frac1M\,\psi_{\mathfrak B_M}+O(1)\ \text{locally on }X.
\end{equation}
Define the ambient current
$Q_M\coloneqq (1-L/M)T_M+(L/M)R\in\alpha$.
For $M>L$, the lower bounds in \eqref{nk:eq:def-TM} and
\eqref{nk:eq:def-R} give
\begin{equation}\label{nk:eq:QM-positivity}
  Q_M\ge\frac{L\delta_R-c_0+Lc_0/M}M\,\omega
  \ge\epsilon_0\,\omega,\qquad
  \epsilon_0\coloneqq \frac{L\delta_R-c_0}M,
\end{equation}
so $Q_M$ is a K\"ahler current, since $L\delta_R>c_0$ by
\eqref{nk:eq:constants}. Its global potential with respect to $\theta$ is
$\varphi_{Q_M}\coloneqq (1-L/M)U_M+(L/M)\Phi_R$;
since the polar set of $R$ is ${E_{\mathrm{nK}}}(\alpha)\supset V$, $\varphi_{Q_M}=-\infty$ at every point of $V$. The
lower bound in \eqref{nk:eq:def-TM} and the estimate \eqref{nk:eq:PhiR-lower}
give, locally on $X$,
\begin{equation}\label{nk:eq:phiQ-lower}
  \varphi_{Q_M}\ge\frac{1-L/M}M\psi_{\mathfrak B_M}
   +\frac{LA_R}M\bigl(\psi_{\mathcal I_V}+\psi_{\mathfrak j}\bigr)+O(1).
\end{equation}

\Needspace{19\baselineskip}
\subsection{The gluing data}

\begin{proposition}\label{nk:prop:principalization}
There exist a sequence of
blowups $\pi\colon Y\to X$ adapted to $V$, with strict transform
$W\subset Y$ of $V$, a point $y\in W$, an open neighborhood $U$
of $W$ in $Y$, a quasi-psh function $\varphi_G$ on $Y$ and a
quasi-psh function $\varphi_U$ on $U$ such that, with
$\theta_Y\coloneqq \pi^*\theta$:
\begin{enumerate}[label=\textup{(\arabic*)},leftmargin=*,itemsep=.35em]
\item\label{nk:it:pr-global} $\theta_Y+\ddc\varphi_G$ is a K\"ahler
  current, and $\varphi_G=-\infty$ at every point of $W$;
\item\label{nk:it:pr-local} $\theta_Y+\ddc\varphi_U$ dominates a K\"ahler
  form on $U$, and $\varphi_U$ is smooth near $y$;
\item\label{nk:it:pr-collar} for every compact set
  $\Lambda\subset U\setminus W$, there is a constant $C_\Lambda$
  such that $\varphi_U\le\varphi_G+C_\Lambda$ on $\Lambda$.
\end{enumerate}
\end{proposition}

The rest of this section is devoted to the proof of
\cref{nk:prop:principalization}. Throughout, when we pull back a divisor
along a further blowup we keep the same letter for its total transform,
and the divisors $G_m$, $G_M$, $G_B$ on the strict transform of $V$
are understood on the current model.

\subsection*{Step 1: the moving ideals}

The modification $\pi$ is built in three stages,
\[
  Y\longrightarrow Y_1\longrightarrow Y_{{\mathrm{fix}}}\longrightarrow X,
\]
with successive strict transforms $W\to W_1\to W_m\to V$ of $V$. The
first stage $\pi_{{\mathrm{fix}}}$, already fixed in
\cref{nk:sec:fixed-principalization}, principalizes $\mathfrak B_m$; the second
principalizes the moving part of $\mathfrak B_M$ relative to
$\mathfrak B_m^{M/m}$; the third principalizes an ideal realizing the excess
divisor $F_B$ of \cref{nk:lem:ideal-comparison}.

Since $\mathfrak B_M\subset\mathfrak B_m^{M/m}$ (\cref{nk:sec:positive-current}),
\eqref{nk:eq:def-Gm-tilde} gives
\begin{equation}\label{nk:eq:def-cM}
  \mathfrak B_M\mathcal O_{Y_{{\mathrm{fix}}}}=\mathcal O_{Y_{{\mathrm{fix}}}}\Bigl(-\frac Mm\widetilde G_m\Bigr)\cdot\mathfrak c_M
\end{equation}
for a coherent ideal $\mathfrak c_M$. Its restriction to $W_m$ is
nonzero because $\overline{\mathfrak B}_M\ne0$. Apply
\cref{nk:lem:relative-principalization} to $\mathfrak c_M$ on $Y_{{\mathrm{fix}}}$, with
$W_m$ as submanifold: we obtain a sequence of blowups
$Y_1\to Y_{{\mathrm{fix}}}$ adapted to $W_m$, with strict transform
$W_1$ of $V$, and an effective exceptional divisor $P^{{\mathrm{mov}}}$ on
$Y_1$ such that $\mathfrak c_M\mathcal O_{Y_1}$ is invertible, and, by
\cref{nk:lem:relative-principalization}\ref{nk:it:rp-restriction},
\begin{equation}\label{nk:eq:Pmov-restriction}
  P^{{\mathrm{mov}}}|_{W_1}=G_M-\frac MmG_m ,
\end{equation}
where $G_M$ denotes the divisor of the invertible ideal
$\overline{\mathfrak B}_M\mathcal O_{W_1}$; indeed, $\NkRes{W_m}\mathfrak c_M
=\overline{\mathfrak B}_M\mathcal O_{W_m}\cdot\mathcal O_{W_m}(\tfrac MmG_m)$ by \eqref{nk:eq:def-cM}. In
particular $\mathfrak B_M$ is invertible on $Y_1$; we write
$\mathfrak B_M\mathcal O_{Y_1}=\mathcal O_{Y_1}(-\widetilde G_M)$, so that
$\widetilde G_M|_{W_1}=G_M$.

Next, by \cref{nk:lem:ideal-comparison}, on $W_1$ the sheaf
\[
  \mathfrak f_B\coloneqq (\mathfrak B\mathcal O_{W_1})\cdot\mathcal O_{W_1}(G_M-pG_m)\subset\mathcal O_{W_1}
\]
is a coherent ideal, and on any modification of $W_1$ on which $\mathfrak B$
is invertible its divisor is $G_B+pG_m-G_M=F_B$. Let
$\widetilde{\mathfrak f}_B\subset\mathcal O_{Y_1}$ be the inverse image of $\mathfrak f_B$
under the surjection $\mathcal O_{Y_1}\to\mathcal O_{W_1}$. This is a coherent ideal
containing $\mathcal I_{W_1}$, with $\NkRes{W_1}\widetilde{\mathfrak f}_B=\mathfrak f_B$, and it
is not contained in $\mathcal I_{W_1}$ because $\mathfrak f_B\ne0$. Apply
\cref{nk:lem:relative-principalization} to $\widetilde{\mathfrak f}_B$ on $Y_1$,
with $W_1$ as submanifold: we obtain a sequence of blowups
$Y\to Y_1$ adapted to $W_1$, with strict transform $W$ of $V$, and
an effective exceptional divisor $P^F$ on $Y$ such that
\begin{equation}\label{nk:eq:PF-restriction}
  P^F|_W=F_B ,
\end{equation}
so that $\mathfrak f_B\mathcal O_W=\mathcal O_W(-F_B)$, and in particular
$\mathfrak B\mathcal O_W=\mathcal O_W(-G_B)$ is invertible on $W$, with
$G_B=F_B+G_M-pG_m$.

Let $\pi\colon Y\to X$ be the composition of $\pi_{{\mathrm{fix}}}$ with the
two sequences of blowups just described. It is a sequence of blowups
adapted to $V$, so $\pi$ is an isomorphism over a nonempty relatively
open subset of $V$, and $W\subset Y$ is the strict transform of $V$.

\subsection*{Step 2: the divisor \texorpdfstring{$H$}{H}}

Let $H\coloneqq H^{{\mathrm{fix}}}_m+P^{{\mathrm{mov}}}+P^F$,
an effective divisor on $Y$ supported on the exceptional locus of
$\pi$, the first two terms being total transforms. Applying
\cref{nk:lem:relative-principalization}\ref{nk:it:rp-ideal-W} to each of the
three sequences of blowups and pulling back, as noted after that lemma,
we obtain
\begin{equation}\label{nk:eq:IV-H}
  \mathcal I_V\mathcal O_Y=\mathcal O_Y(-H)\cdot\mathcal I_W .
\end{equation}

\subsection*{Step 3: the K\"ahler current \texorpdfstring{$\theta_Y+\ddc\varphi_G$}{on Y}}

Apply \cref{nk:lem:exceptional-metric} to $\pi$ and $\omega$: there is
an effective $\pi$-exceptional $\mathbf Q$-divisor $E_{{\mathrm{rel}}}$, with
a smooth metric, such that $\omega_Y\coloneqq \pi^*\omega-\Theta_{E_{{\mathrm{rel}}}}$
is a K\"ahler form on $Y$. Since $\pi$ is an isomorphism over a
nonempty relatively open subset of $V$, no component of $E_{{\mathrm{rel}}}$
contains $W$. Fix a rational number $\tau$ with
$0<\tau<\min\{\epsilon_0,\delta_0/2\}$, and set
$\varphi_G\coloneqq \varphi_{Q_M}\circ\pi+\tau\psi_{E_{{\mathrm{rel}}}}$. By
\eqref{nk:eq:QM-positivity},
\[
  \theta_Y+\ddc\varphi_G
  =\pi^*Q_M+\tau\bigl([E_{{\mathrm{rel}}}]-\Theta_{E_{{\mathrm{rel}}}}\bigr)
  \ge(\epsilon_0-\tau)\pi^*\omega+\tau\omega_Y
  \ge\tau\omega_Y ,
\]
so $\theta_Y+\ddc\varphi_G$ is a K\"ahler current. Since
$\varphi_{Q_M}=-\infty$ on $V$, $\varphi_G=-\infty$ at every point of $W$. This proves
\ref{nk:it:pr-global}.

\subsection*{Step 4: the divisor \texorpdfstring{$D$}{D} and a K\"ahler class on \texorpdfstring{$W$}{W}}

Define
\begin{equation}\label{nk:eq:def-DLY}
  D\coloneqq \frac1M\widetilde G_M+\frac{C-1}MP^{{\mathrm{mov}}}+\frac CMP^F+\tau E_{{\mathrm{rel}}}.
\end{equation}
This is an effective
$\mathbf Q$-divisor on $Y$, and no component of it
contains $W$. Since total transforms commute
with restriction to strict transforms, \eqref{nk:eq:GB-decomposition},
\eqref{nk:eq:Pmov-restriction} and \eqref{nk:eq:PF-restriction} give
\begin{equation}\label{nk:eq:DL-restriction}
  D|_W=\frac1MG_M+\frac{C-1}M\Bigl(G_M-\frac MmG_m\Bigr)+\frac CMF_B
       +\tau E_{{\mathrm{rel}}}|_W
  =\frac1\ell G_B+\tau E_{{\mathrm{rel}}}|_W .
\end{equation}

The restriction $\pi|_W\colon W\to V$ is a modification of $V$.
Since $\varphi_{K,\ell}$ has analytic singularities of type
$(\ell^{-1},\mathfrak B)$ and $\mathfrak B\mathcal O_W=\mathcal O_W(-G_B)$ is invertible, the
function
$\varphi'\coloneqq (\pi|_W)^*\varphi_{K,\ell}+\tau\psi_{E_{{\mathrm{rel}}}}|_W$
has log singularities along $D|_W$ with the restricted metric, by
\eqref{nk:eq:DL-restriction}. Put
$\theta'\coloneqq (\pi^*\theta-\tfrac12\delta_0\pi^*\omega)|_W$. Since
$(\pi|_W)^*K_\ell\ge\delta_0\,\pi^*\omega|_W$ by \eqref{nk:eq:def-K-ell},
\[
 \begin{aligned}
  \theta'+\ddc\varphi'
  &=(\pi|_W)^*K_\ell-\tfrac12\delta_0\pi^*\omega|_W
   +\tau\bigl([E_{{\mathrm{rel}}}|_W]-\Theta_{E_{{\mathrm{rel}}}}|_W\bigr)\\
  &\ge(\tfrac12\delta_0-\tau)\pi^*\omega|_W+\tau\omega_Y|_W
  \ge\tau\omega_Y|_W ,
 \end{aligned}
\]
so $\theta'+\ddc\varphi'$ is a K\"ahler current on $W$. By
\cref{nk:lem:residual-nef}, the class
\begin{equation}\label{nk:eq:AW-class}
  \{\theta'\}-\{D|_W\}
  =\{\theta_Y-\Theta_D-\tfrac12\delta_0\pi^*\omega\}|_W
\end{equation}
is K\"ahler on $W$.

\subsection*{Step 5: the correction term and the comparison}

Set $\chi\coloneqq \frac{C-1}M\,\Psi_m\circ\pi$,
with $\Psi_m$ as in \eqref{nk:eq:def-Psi}. Then $\chi\le0$ is quasi-psh
on $Y$, and it is smooth on the open set
$\pi^{-1}(X\setminus(Z'\cup Z(\mathfrak B_m)))$. This open set meets
$W$, since its complement on $V$, together with the locus over which
$\pi$ is not an isomorphism, is a proper analytic subset of $V$.
Hence $\chi$ is smooth on a neighborhood of some point of $W$. Moreover
$\ddc\chi\ge-((C-1)c_\Psi/M)\pi^*\omega$, which is at least
$-\tfrac12\delta_0\pi^*\omega$ because
$M\ge2(C-1)c_\Psi/\delta_0$ by the choice of $M$ in
\cref{nk:sec:positive-current}.

We claim that for every compact set $\Lambda\subset Y\setminus W$
there is a constant $C_\Lambda$ such that
\begin{equation}\label{nk:eq:collar}
  \psi_D+\chi\le\varphi_G+C_\Lambda\qquad\text{on }\Lambda.
\end{equation}
By \cref{nk:sec:qpsh}, locally on $Y$ we have
$\psi_{\mathfrak B_M}\circ\pi=\psi_{\widetilde G_M}+O(1)$,
$\psi_{\mathfrak B_m}\circ\pi=\psi_{\widetilde G_m}+O(1)$ and, by
\eqref{nk:eq:IV-H},
$\psi_{\mathcal I_V}\circ\pi=\psi_H+\psi_{\mathcal I_W}+O(1)$, where
$\psi_{\mathcal I_W}$ is bounded on $\Lambda$. Hence \eqref{nk:eq:phiQ-lower} gives
\begin{equation}\label{nk:eq:QM-lower}
  \varphi_{Q_M}\circ\pi\ge\frac{1-L/M}M\psi_{\widetilde G_M}
   +\frac{LA_R}M\bigl(\psi_H+\psi_{\mathfrak j}\circ\pi\bigr)-C_\Lambda
  \qquad\text{on }\Lambda
\end{equation}
for some constant $C_\Lambda$. On the other hand, by \eqref{nk:eq:def-Psi} and
\eqref{nk:eq:def-Gm-tilde},
\[
  \chi=\frac{C-1}M\bigl(\psi_{\mathfrak j}\circ\pi+\psi_{\widetilde G_m}\bigr)+O(1)
  \le\frac{C-1}M\bigl(\psi_{\mathfrak j}\circ\pi+\psi_{H^{{\mathrm{fix}}}_m}\bigr)+O(1)
  \qquad\text{on }Y,
\]
and $\psi_H=\psi_{H^{{\mathrm{fix}}}_m}+\psi_{P^{{\mathrm{mov}}}}+\psi_{P^F}+O(1)$ by
the definition of $H$. By \eqref{nk:eq:constants},
$C-1\ge LA_R$ and hence $C\ge LA_R$, while $1\ge1-L/M$. Since all
the functions
$\psi_{\mathfrak j}\circ\pi$,
$\psi_{\widetilde G_M}$, $\psi_{H^{{\mathrm{fix}}}_m}$, $\psi_{P^{{\mathrm{mov}}}}$,
$\psi_{P^F}$ are bounded above, we conclude from \eqref{nk:eq:def-DLY} that
\[
  \psi_{D-\tau E_{{\mathrm{rel}}}}+\chi
  \le\frac{1-L/M}M\psi_{\widetilde G_M}
   +\frac{LA_R}M\bigl(\psi_H+\psi_{\mathfrak j}\circ\pi\bigr)+O(1)
  \qquad\text{on }Y.
\]
Adding $\tau\psi_{E_{{\mathrm{rel}}}}$ to both sides and using
\eqref{nk:eq:QM-lower} and the definition of $\varphi_G$, we obtain
\eqref{nk:eq:collar}.

\subsection*{Step 6: the local current \texorpdfstring{$\theta_Y+\ddc\varphi_U$}{near W}}

Choose a K\"ahler form $A_W$ in the class \eqref{nk:eq:AW-class}, and a smooth function $f$ on $W$ such
that
$A_W=(\theta_Y-\Theta_D-\tfrac12\delta_0\pi^*\omega)|_W+\ddc f$.
Extend $f$ to a smooth function $\widetilde f$ near $W$, and let
$r_W^2$ be the square of the distance to $W$ with respect to a
Riemannian metric on $Y$. For $B\gg1$, there is a relatively compact
neighborhood $U$ of $W$ on a neighborhood of whose closure
\[
  \Omega_U\coloneqq \theta_Y-\Theta_D-\tfrac12\delta_0\pi^*\omega
       +\ddc(\widetilde f+Br_W^2)>0 :
\]
indeed, along $W$ the restriction of $\Omega_U$ to $T_W$ is
$A_W$, while the Levi form of $r_W^2$ is positive semidefinite with
kernel exactly $T_W$. In particular, $\Omega_U$ dominates a K\"ahler
form on $U$. Set
$\varphi_U\coloneqq \widetilde f+Br_W^2+\psi_D+\chi$ on $U$. By Step~5,
\[
  \theta_Y+\ddc\varphi_U
  =\Omega_U+[D]+(\tfrac12\delta_0\pi^*\omega+\ddc\chi)
  \ge\Omega_U ,
\]
so $\theta_Y+\ddc\varphi_U$ dominates a K\"ahler form on $U$.
Since no component of $D$ contains $W$, the set $\Supp D\cap W$
is a proper analytic subset of $W$; as $\chi$ is smooth near some
point of $W$, we may choose $y\in W$ outside $\Supp D$ near which
$\chi$ is smooth. Then $\varphi_U$ is smooth near $y$, which
proves \ref{nk:it:pr-local}. Finally, if $\Lambda\subset U\setminus W$ is
compact, then $\widetilde f+Br_W^2$ is bounded on $\Lambda$, and
\eqref{nk:eq:collar} gives $\varphi_U\le\varphi_G+C'_\Lambda$ on
$\Lambda$ for a constant $C'_\Lambda$, which is
\ref{nk:it:pr-collar}. This completes the proof of
\cref{nk:prop:principalization}.\qed

\subsection{Completion of the proof}\label{nk:sec:gluing}

Let $\pi\colon Y\to X$, $W$, $y$, $U$, $\varphi_G$ and
$\varphi_U$ be as in \cref{nk:prop:principalization}, and choose open
sets $W\Subset U_1\Subset U_2\Subset U$. The hypotheses of
\cref{nk:prop:gluing} are satisfied for $\theta_Y=\pi^*\theta$,
$\varphi_G$, $\varphi_U$, $U_1$, $U_2$ and the point
$y\in W\subset U_1$: \ref{nk:it:gl-current} and \ref{nk:it:gl-local} are
\cref{nk:prop:principalization}\ref{nk:it:pr-global} and
\ref{nk:it:pr-local}, and \ref{nk:it:gl-collar} is
\cref{nk:prop:principalization}\ref{nk:it:pr-collar} for the compact set
$\Lambda=\overline U_2\setminus U_1$. We obtain a K\"ahler current in
$\pi^*\alpha$ which is smooth near $y$, and so
$y\notin{E_{\mathrm{nK}}}(\pi^*\alpha)$. On the other hand, the transformation rule
\eqref{nk:eq:transformation-rules} gives
\[
  {E_{\mathrm{nK}}}(\pi^*\alpha)=\pi^{-1}\bigl({E_{\mathrm{nK}}}(\alpha)\bigr)\cup\operatorname{Exc}(\pi)
  \supset\pi^{-1}(V)\supset W\ni y,
\]
a contradiction. Therefore the assumption $h>0$ made in
\cref{nk:sec:setup} is false, and
$\langle\alpha^k\rangle_{X\mid V}=0$.
This proves the first assertion of \cref{nk:thm:main}.

For the second assertion, if $\alpha$ is not big, then
${E_{\mathrm{nK}}}(\alpha)=X$ by convention, and $X\subset\operatorname{Null}(\alpha)$ because
$\langle\alpha^n\rangle_{X\mid X}=\vol(\alpha)=0$; hence
$\operatorname{Null}(\alpha)=X={E_{\mathrm{nK}}}(\alpha)$. If $\alpha$ is big, then
${E_{\mathrm{nK}}}(\alpha)$ has no isolated points by \cref{prop:bc-no-isolated-locus}, so
each of its irreducible components is positive-dimensional and, by the
first assertion, contained in $\operatorname{Null}(\alpha)$; thus
${E_{\mathrm{nK}}}(\alpha)\subset\operatorname{Null}(\alpha)$. Conversely, every
positive-dimensional irreducible subvariety $V\not\subset{E_{\mathrm{nK}}}(\alpha)$
has positive numerical restricted volume by
\cref{nk:sec:restricted-volumes}\ref{nk:it:rv-positive}, so every subvariety
occurring in the union \eqref{nk:eq:null-locus} is contained in
${E_{\mathrm{nK}}}(\alpha)$, and $\operatorname{Null}(\alpha)\subset{E_{\mathrm{nK}}}(\alpha)$. This completes
the proof of \cref{nk:thm:main}.\qed


\section{What carries over to a Hermitian manifold}

The proof separates local analytic tools from numerical consequences of
closedness. The regularization lemma \ref{nk:lem:regularization} and the
global ideal potential lemma \ref{nk:lem:ideal-function} apply on every
compact Hermitian manifold. In the first lemma, the reference metric
measures the lower bound; the background $\zeta$ is still closed. Local
potentials for $\zeta$ give holomorphic changes of weight on overlaps,
while Euclidean volume measures allow the local $L^2$ estimates to be
used without a global K\"ahler form. No Monge--Amp\`ere mass bound enters
this construction.

The divisorial domination argument uses more. Its residual classes are
nef on K\"ahler resolutions, their mixed intersections are controlled on
the fixed base, and a numerical inequality produces a big class. The
trace formula, the mass continuity in $d_S$ and the rooftop criterion are
also used in their K\"ahler regimes. The local regularization theorem
does not supply substitutes for these steps on a general Hermitian
manifold.

There is a further difficulty for numerical restricted volume itself.
The unrestricted Hermitian-error supremum is infinite in the zero class
on the example of \cref{prop:bc-zero-class-infinite-error-volume}.
Thus a Hermitian analogue of the argument must first specify a finite
restricted invariant and prove its relation to an appropriate trace.
The qualitative non-Hermitian locus of a nonclosed big background,
discussed in \cref{sec:non-kahler-geometric-problems}, addresses a
different positivity condition. The present theorem establishes the
closed K\"ahler case in full; it identifies precise additional steps
needed for a broader theory.
